\chapter{Cells: structure and function}
\label{vol06:ch01}

\epigraph{Living matter, while not eluding the ``laws of physics'' as
established up to date, is likely to involve ``other laws of physics''
hitherto unknown, which, however, once they have been revealed, will
form just as integral a part of this science as the former.}{Erwin
Schr\"odinger, \emph{What Is Life?} (1944), ch.~7, in the Cambridge
edition \cite[p.~76]{hist:schrodinger1944}.}

\chapmeta{Half-life: foundational. Archetypes: interface, control. Exercise target: 25. Project track: household microscopy and observation. Hazard class: low (visual inspection of pond water, yeast, and yoghurt cultures with a low-cost optical microscope). Reader path default: standard, with sections explicitly tagged. Named cases: none required in this chapter; cellular failure modes are described as mechanism classes (apoptosis, necrosis, senescence).}

\noindent
Volume VI opens with the cell because every other claim the volume will
make about living matter is, in the end, a claim about an assembly of
cells. The cell is the smallest unit of life that satisfies the
engineering definition: a bounded region of matter that maintains its
own state against thermodynamic disorder, that processes information,
that replicates with controlled error, and that fails by named
mechanisms. The chapter reads the cell as such a bounded system. It
identifies the components, names the control loops, walks through four
worked cellular archetypes, and closes by naming the three principal
modes of cellular failure that will recur throughout the volume.

\section{The cell as engineered object}\pathtag{core}

A \engterm{cell} is a self-bounded, self-maintaining unit of living
matter. The boundary is a phospholipid bilayer membrane. The interior
is an aqueous compartment in which proteins, nucleic acids, sugars,
and lipids are present at concentrations that lie far outside any
thermodynamic equilibrium with the cell's surroundings. The departure
from equilibrium is the cell's defining property; the membrane is the
device that maintains it.

The departure is large. A typical mammalian cell contains potassium
ions at a concentration of $140\,\text{mM}$ while sitting in extracellular
fluid at $4\,\text{mM}$; sodium ions are $10\,\text{mM}$ inside and
$145\,\text{mM}$ outside; calcium ions are $0.0001\,\text{mM}$ inside
and $1\,\text{mM}$ outside \cite[ch.~11]{text:v6c01-alberts-mbc-2022}.
The ratios are factors of $35$, $14$, and $10\,000$ respectively.
Each ratio is maintained against a passive gradient that would
otherwise drive the cell to ionic uniformity within seconds. The cell
is doing work continuously to maintain the disequilibrium, and the
work is the metabolic work that Chapter~\ref{vol06:ch02} will
quantify.

The cell's typical size sets the engineering scale of the volume. A
bacterial cell is $1$ to $2\,\si{\micro\meter}$ across and contains
about $4 \times 10^{-15}\,\text{kg}$ of dry material. A typical
animal cell is $10$ to $30\,\si{\micro\meter}$ across and contains
about $10^{-12}\,\text{kg}$, with the dry mass roughly half of the
total. A red blood cell is $7.5\,\si{\micro\meter}$ across by
$2\,\si{\micro\meter}$ thick. A mature human ovum is
$120\,\si{\micro\meter}$, the largest cell visible to the unaided eye.
The number of cells in an adult human body is approximately
$3.7 \times 10^{13}$, of which roughly $2.5 \times 10^{13}$ are red
blood cells \cite{paper:v6c01-bianconi-2013}. The bacterial cells
associated with a human host (the microbiota) number roughly the same
order, with current estimates putting the bacterial-to-human cell
ratio at approximately $1.3:1$ rather than the older $10:1$ figure
\cite{paper:v6c01-sender-2016}.
\Cref{fig:vol06:ch01:size-scaling} plots characteristic cell sizes
on a log-log axis; the range spans five orders of magnitude in linear
size and twelve in volume.

% Figure: log-log plot of cell type vs. characteristic linear size.
% Schmidt-Nielsen style. Pure tikz so no external data file needed.
\begin{figure}[htbp]
\centering
\begin{tikzpicture}[
  axis/.style={thick, draw=engblack, -{Latex[length=2.0mm]}},
  pt/.style={circle, fill=engblue, draw=engblue!80!black, inner sep=1.6pt},
  lbl/.style={font=\scriptsize, anchor=west, align=left},
  band/.style={fill=engamber!10, draw=none},
]

% Axes from log10(L/m) = -7 to -2 (100 nm to 10 mm)
% x range: -7 (100 nm) to -2 (10 mm); width 11 cm => 2.2 cm per decade
% y range: -7 to 7 for log10(volume / um^3); height 8 cm => 0.57 cm per unit

\draw[axis] (0,0) -- (11.5, 0);
\draw[axis] (0,0) -- (0, 8.2);
\node[font=\small] at (5.75, -0.7) {Characteristic linear size $L$};
\node[font=\small, rotate=90, anchor=south] at (-0.9, 4.0)
  {Volume scaling $L^3$ ($\mu$m$^3$)};

% x-axis tick labels (decades)
\foreach \d/\lbl in {0/{100 nm}, 2.2/{1 \,$\mu$m}, 4.4/{10 \,$\mu$m},
                     6.6/{100 \,$\mu$m}, 8.8/{1 mm}, 11.0/{10 mm}} {
  \draw[thick] (\d, 0) -- (\d, -0.12);
  \node[anchor=north, font=\scriptsize] at (\d, -0.12) {\lbl};
}
% y-axis tick labels (log10 volume in um^3)
\foreach \d/\lbl in {0/{$10^{-3}$}, 1.14/{$10^{-2}$}, 2.28/{$10^{-1}$},
                     3.42/{$1$}, 4.56/{$10^{1}$}, 5.7/{$10^{3}$},
                     6.84/{$10^{6}$}, 7.98/{$10^{9}$}} {
  \draw[thick] (-0.12, \d) -- (0, \d);
  \node[anchor=east, font=\scriptsize] at (-0.15, \d) {\lbl};
}

% Reference line: V = L^3 (so log V = 3 log L). For L in m at decade d
% from -7, log10(L/um) = d/2.2 - 1, and log10(V/um^3) = 3 * (d/2.2 - 1).
% That maps to y = 3*(d/2.2 - 1) * 1.14 in unit-multiplied space.
% Cleanly: y_per_x ratio = 3 * 1.14 / 2.2 = 1.554.
\draw[densely dashed, draw=enggray, thick]
  (0, 0) -- (11.0, 7.7);
\node[anchor=west, font=\scriptsize\itshape] at (8.7, 7.7) {$V = L^3$};

% Data points (approximate; canonical cell-size compilation)
% Format: (decade_x, decade_y, label, label_y_offset)
% Decade_x: position from -7 in cm (2.2 per decade)
% Decade_y: log10(V/um^3) * 1.14

% Mycoplasma:    L ~ 200 nm  => x ~ 0.66; V ~ 0.004 um^3 => y ~ -0.4 cm (off-axis)
\node[pt] at (0.66, 0.05) {};
\node[lbl] at (0.95, 0.0) {Mycoplasma};

% E. coli:       L ~ 1 um    => x = 2.2; V ~ 1 um^3 => y = 3.42
\node[pt] at (2.20, 3.42) {};
\node[lbl] at (2.50, 3.42) {\emph{E.\ coli}};

% Yeast:         L ~ 5 um    => x ~ 3.74; V ~ 50 um^3 => y ~ 5.0 (10^1.7)
\node[pt] at (3.74, 4.50) {};
\node[lbl] at (4.05, 4.50) {Yeast};

% Erythrocyte:   L ~ 8 um    => x ~ 4.25; V ~ 90 um^3 => y ~ 5.16
\node[pt] at (4.25, 5.20) {};
\node[lbl] at (4.55, 5.18) {RBC};

% Lymphocyte:    L ~ 10 um   => x = 4.4; V ~ 200 um^3 => y ~ 5.27
\node[pt] at (4.40, 5.40) {};
\node[lbl] at (4.70, 5.50) {Lymphocyte};

% Fibroblast:    L ~ 30 um   => x ~ 5.45; V ~ 2000 um^3 => y ~ 5.7
\node[pt] at (5.45, 5.70) {};
\node[lbl] at (5.75, 5.70) {Fibroblast};

% Cardiomyocyte: L ~ 100 um (length)  => x = 6.6; V ~ 20000 um^3 => y ~ 6.30
\node[pt] at (6.60, 6.30) {};
\node[lbl] at (6.90, 6.30) {Cardiomyocyte};

% Adipocyte:     L ~ 100 um  => x = 6.6; V ~ 500000 um^3 => y ~ 6.84
\node[pt] at (6.60, 6.84) {};
\node[lbl] at (6.90, 6.84) {Adipocyte};

% Ovum:          L ~ 120 um  => x ~ 6.78; V ~ 9e5 um^3 => y ~ 7.13
\node[pt] at (6.78, 7.13) {};
\node[lbl] at (7.05, 7.13) {Ovum};

% Neuron axon:   L ~ 1 m (length); but width ~ 10 um. Annotate special.
\node[pt, fill=engred] at (10.5, 5.50) {};
\node[lbl] at (10.0, 6.05) {Motor neuron \\ (axon length)};

\end{tikzpicture}
\caption[Cell-size scaling]{Characteristic linear size and approximate
volume of representative cell types, plotted on a log-log axis. The
range spans five orders of magnitude in linear size and roughly twelve
in volume; the dashed reference line is $V = L^3$. Bacterial cells
sit near $1\,\si{\micro\meter}$; typical animal cells at $10$ to
$30\,\si{\micro\meter}$; specialised cells (ovum, cardiomyocyte)
push the volume upward without leaving the line. The motor neuron is
the outlier: its axon length can exceed a metre while its diameter
stays at the typical-cell scale, so its linear size depends on which
dimension is measured. Compiled from canonical values
\cite{text:v6c01-bionumbers-2015,paper:v6c01-bianconi-2013}.}
\label{fig:vol06:ch01:size-scaling}
\end{figure}


\engterm{The word ``cell.''} The English word ``cell,'' applied to a
biological structure, dates to Robert Hooke's 1665 \emph{Micrographia},
observation~XVIII, where Hooke described the prismatic compartments
he saw in thin slices of cork as ``little boxes or cells distinguished
from one another'' \cite[obs.~XVIII]{hist:v6c01-hooke-1665}. Hooke was
looking at the dead, lignified cell walls of cork tissue; the term
was extended to the living units of plant and animal tissue only in
the nineteenth century, after Schleiden and Schwann's cell theory
crystallised the claim that all living matter is built of cells. The
etymology is worth recording: ``cell'' began as a name for a
\emph{compartment} and is still, in the engineering reading, primarily
a name for a bounded compartment that does work against its
surroundings.

\engterm{Reading the cell as an engineered object}. The engineer who
opens a service manual for an unfamiliar machine asks four questions
in sequence: what does it do, what does it use as input, what does it
produce as output, and what controls the rate. The same four questions
organise a productive reading of a cell. A red blood cell takes oxygen
from the lung capillary, binds it to haemoglobin, releases it at
distant capillaries, and returns. The input is dissolved oxygen and
glucose; the output is delivered oxygen and dissolved carbon dioxide;
the controller is the local oxygen partial pressure and the local pH.
The cell is a transport device with a state-dependent affinity curve;
the affinity curve is the sigmoidal oxygen-haemoglobin saturation
curve, and its sigmoidicity is the cell's signature engineering feature.

We adopt that reading discipline throughout the volume. A cell is
named by its function, its inputs, its outputs, and the regulators of
its rate. Histological taxonomy (epithelial, muscle, nerve, connective)
is downstream of the engineering taxonomy and is introduced where it
clarifies a mechanism.

\engterm{The cell as a chemical reactor}. The cell is also, in the
chemical-engineering sense, a stirred-tank reactor: an aqueous
compartment in which a few thousand simultaneous reactions run at
rates set by enzyme concentrations and substrate availability. The
fluxes through the network sum to a metabolic accounting that
Chapter~\ref{vol06:ch02} will balance. The chemical-reactor reading
clarifies what the cell is doing thermodynamically; the
information-processing reading, taken up in Chapter~\ref{vol06:ch03}
and Chapter~\ref{vol06:ch04}, clarifies what it is doing
algorithmically. The two readings are complementary; the volume's
discipline is to keep both in view.

\engterm{What the cell is not}. The cell is not a static object. Every
component in the cell is being synthesised, degraded, and replaced on
a timescale set by its function: structural proteins turn over on a
timescale of days, signalling proteins on a timescale of minutes,
metabolic intermediates on a timescale of seconds. A cell that stops
turning over its components dies within hours from the accumulated
damage that turnover normally removes
\cite[ch.~6]{text:v6c01-alberts-mbc-2022}. The cell is not a
fixed-composition object either: the same cell at different points in
the cell cycle, or in different metabolic states, has different
concentrations of the same molecules by factors of ten or more. The
cell is a dynamic equilibrium, and the engineering reading treats it
as such.

\engterm{Concentrations at engineering working depth}. A typical
mammalian cytoplasm carries protein at approximately $200$ to
$300\,\text{mg}/\text{mL}$, or roughly $3$ to $5\,\text{mM}$ averaged
over an average protein of $50\,\text{kDa}$
\cite{paper:v6c01-luby-phelps-2000}. The protein crowding reduces
small-molecule diffusion coefficients by a factor of $3$ to $4$
relative to dilute solution, and protein diffusion coefficients by an
order of magnitude or more. Cytoplasmic water is approximately
$70\,\%$ of cell mass; ATP is $1$ to $5\,\text{mM}$; ribosomes are
present at approximately $1$ ribosome per $10^{-2}\,\si{\micro\meter}^3$
in a typical mammalian cell, or some $4$ million ribosomes total. The
canonical numbers compendium is the \emph{Cell Biology by the
Numbers} project, which curates the BioNumbers database
\cite{text:v6c01-bionumbers-2015,paper:v6c01-bionumbers-database-2010};
we cite it throughout the volume as the working reference for
order-of-magnitude reasoning.

\begin{principle}[A cell is a bounded, maintained, disequilibrium system]
A cell maintains itself far from thermodynamic equilibrium with its
surroundings by continuous expenditure of free energy across a
semi-permeable boundary. The disequilibrium is what makes the cell
alive; the boundary is what makes it a cell; the energy expenditure is
what links it to Chapter~\ref{vol06:ch02}. A cell that loses any one
of the three (membrane integrity, free-energy supply, or the
controlled disequilibrium itself) ceases to be a cell on a timescale
of seconds to hours, depending on which is lost.
\end{principle}

\section{Membranes, organelles, cytoskeleton}\pathtag{core}

The internal architecture of a cell is built from three classes of
component: membranes (the boundary itself and the membranes of internal
compartments), organelles (membrane-bounded subcompartments with
specialised function), and the cytoskeleton (the protein-fibre network
that gives the cell its mechanical structure and that organises the
transport of organelles and vesicles).

\engterm{The plasma membrane}. The cell membrane is a bilayer of
phospholipid molecules, $4$ to $5\,\si{\nano\meter}$ thick. Each
phospholipid has a hydrophilic phosphate head group and two
hydrophobic fatty-acid tails. In water the molecules self-assemble
into the bilayer because the hydrophobic tails minimise their contact
with water by facing each other in the membrane core, while the
hydrophilic heads face the aqueous interior and exterior. The
self-assembly is spontaneous; the resulting bilayer is fluid (the
phospholipid molecules diffuse laterally within the leaflet on a
timescale of microseconds per molecular diameter) and self-healing
(small punctures reseal because the hydrophobic tails are unstable
when exposed to water). The canonical model of the membrane as a
two-dimensional viscous fluid with embedded proteins is Singer and
Nicolson's 1972 fluid-mosaic model
\cite{paper:v6c01-singer-nicolson-1972}, which remains the textbook
picture and the basis of every subsequent refinement.

The bilayer is decorated with proteins that perform the membrane's
selective functions. Three structural classes are useful to name:
\engterm{integral membrane proteins} that span the bilayer (channels,
pumps, receptors); \engterm{peripheral membrane proteins} that
associate with one face of the bilayer (signalling adaptors,
cytoskeletal anchors); and \engterm{lipid-anchored proteins} that
carry a covalently attached lipid that inserts into one leaflet (the
Ras family of signalling proteins is the textbook example). The
density of proteins in a typical eukaryotic plasma membrane is
$10^3$ to $10^4$ per square micrometre; the proteins occupy roughly
half of the membrane area in many cell types
\cite[ch.~10]{text:v6c01-alberts-mbc-2022}.
\Cref{fig:vol06:ch01:membrane-bilayer} draws the bilayer with its
three protein classes labelled.

% Figure: lipid bilayer with embedded proteins. Phospholipid heads as
% filled circles, tails as wavy lines. Integral, peripheral, and
% lipid-anchored proteins drawn schematically.
\begin{figure}[htbp]
\centering
\begin{tikzpicture}[
  head/.style={circle, fill=engblue, draw=engblue, inner sep=1.4pt},
  tail/.style={engblack, very thick},
  prot/.style={fill=engamber!40, draw=engamber!80!black, very thick},
  channel/.style={fill=engred!30, draw=engred!80!black, very thick},
  lbl/.style={font=\scriptsize, align=center},
]
% Outer (top) leaflet heads
\foreach \x in {0,0.45,...,12.0} {
  \node[head] at (\x, 2.0) {};
  \draw[tail] (\x, 1.95) -- ++(0.05, -0.85) -- ++(-0.05, -0.05);
}
% Inner (bottom) leaflet heads
\foreach \x in {0,0.45,...,12.0} {
  \node[head] at (\x, 0.0) {};
  \draw[tail] (\x, 0.05) -- ++(0.05, 0.85) -- ++(-0.05, 0.05);
}

% Integral channel protein (spanning both leaflets, with a pore)
\draw[channel, rounded corners=2pt]
  (2.6, -0.2) -- (3.4, -0.2) -- (3.4, 2.2) -- (2.6, 2.2) -- cycle;
\draw[white, fill=white] (2.9, 0.3) rectangle (3.1, 1.7);
\draw[channel, very thick] (2.9, 0.3) -- (2.9, 1.7);
\draw[channel, very thick] (3.1, 0.3) -- (3.1, 1.7);

% Integral pump protein (asymmetric, spanning bilayer)
\draw[prot, rounded corners=2pt]
  (5.6, -0.3) -- (6.6, -0.3) -- (6.7, 2.3) -- (5.5, 2.3) -- cycle;
\node[font=\tiny] at (6.1, 1.0) {ATP};

% Peripheral protein (on inner leaflet, doesn't cross)
\draw[prot, rounded corners=2pt]
  (8.4, -0.6) -- (9.6, -0.6) -- (9.6, -0.05) -- (8.4, -0.05) -- cycle;

% Lipid-anchored protein (on outer leaflet)
\draw[prot, rounded corners=2pt]
  (10.3, 2.3) -- (11.4, 2.3) -- (11.4, 2.85) -- (10.3, 2.85) -- cycle;
\draw[tail, thick] (10.6, 2.3) -- ++(0.05, -0.55) -- ++(-0.05, -0.05);
\draw[tail, thick] (10.8, 2.3) -- ++(0.05, -0.55) -- ++(-0.05, -0.05);

% Labels
\node[lbl, anchor=west] at (12.4, 2.4) {extracellular};
\node[lbl, anchor=west] at (12.4, -0.4) {cytoplasm};

\node[lbl, anchor=north, font=\scriptsize] at (3.0, -0.4) {channel};
\node[lbl, anchor=north, font=\scriptsize] at (6.1, -0.4) {pump};
\node[lbl, anchor=south, font=\scriptsize] at (9.0, -0.85) {peripheral};
\node[lbl, anchor=south, font=\scriptsize] at (10.85, 2.9) {lipid-anchored};

% Bilayer thickness label
\draw[<->, thick] (-0.5, 0) -- (-0.5, 2.0)
  node[midway, left, font=\scriptsize] {$\approx 4$ nm};

\end{tikzpicture}
\caption[The plasma membrane: bilayer with embedded proteins]{The
plasma membrane as a phospholipid bilayer with three classes of
embedded protein: integral membrane proteins (channels and pumps)
that span both leaflets; peripheral membrane proteins that associate
with one face of the bilayer without crossing; and lipid-anchored
proteins that carry a covalently attached lipid that inserts into one
leaflet. The bilayer is approximately $4$ nm thick. Proteins occupy
roughly half of the membrane area in many cell types
\cite[ch.~10]{text:v6c01-alberts-mbc-2022}.}
\label{fig:vol06:ch01:membrane-bilayer}
\end{figure}


\engterm{The lipid bilayer as a capacitor}. The bilayer is electrically
a thin dielectric layer separating two conducting aqueous compartments.
The dielectric constant of the hydrophobic core is approximately
$\varepsilon_r = 2$ to $4$; the membrane capacitance per unit area is
$C_m = \varepsilon_0 \varepsilon_r/d \approx 0.9\,\si{\micro\farad}/
\text{cm}^2$, with $d \approx 4\,\si{\nano\meter}$. The membrane is a
capacitor with one of the largest specific capacitances in the
engineered or natural world. The fact has consequences across the
volume: it sets the time constant for the changes in membrane voltage
that propagate signals in neurons (Chapter~\ref{vol06:ch07}); it sets
the charge required to depolarise the membrane by a given amount; and
it sets the dielectric breakdown limit at which electroporation
becomes the mechanism of membrane disruption.

\engterm{The bending elasticity of the bilayer}. The phospholipid
bilayer behaves mechanically like a thin elastic sheet with bending
modulus $\kappa \approx 10\,k_B T \approx 4 \times 10^{-20}\,\text{J}$
\cite{paper:v6c01-helfrich-1973,text:v6c01-phillips-physical-2012}.
The energy of bending the membrane into a closed sphere of any
radius is $8\pi \kappa \approx 1000\,k_B T$, independent of the
sphere's radius; the energy of bending into a tube of length $L$ and
radius $R$ is $\pi\kappa L / R$, which favours large radii. The
bending energy is what sets the shapes of organelles, the formation
of vesicles by clathrin-coated pits, and the deformability of red
blood cells through capillaries; we return to it in
Chapter~\ref{vol06:ch06}.

\engterm{Organelles}. The eukaryotic cell partitions its interior into
membrane-bounded compartments, each with its own internal chemistry.
The principal organelles are the nucleus (carrying the cell's DNA,
bounded by a double membrane perforated by nuclear pores); the
endoplasmic reticulum (ER, an extensive system of folded membranes
where most protein synthesis and lipid synthesis occur); the Golgi
apparatus (a stack of membrane discs that processes and sorts proteins
exiting the ER); mitochondria (double-membrane organelles that
generate most of the cell's ATP by oxidative phosphorylation,
Chapter~\ref{vol06:ch02}); lysosomes (acidic compartments that degrade
imported molecules and damaged organelles); peroxisomes (small
oxidative compartments); and, in plants, chloroplasts (the
photosynthetic organelle) and the central vacuole (a large
osmotic-pressure organ).

The mitochondrion is the cell's principal energy converter. Mitochondrial
DNA, present in $1000$ to $10\,000$ copies per cell, encodes a small
subset of mitochondrial proteins; the bulk of the mitochondrial
proteome is encoded in the nucleus and imported. Mitochondrial
biogenesis is regulated by the cell's energy demand, and the
mitochondrial content of a cell scales with the cell's oxygen
consumption: a cardiomyocyte's volume is $30\,\%$ to $40\,\%$
mitochondria; a fibroblast's, $5\,\%$. The mitochondrial scaling will
appear repeatedly when we account for the cell's energy budget. The
mitochondrion's origin is endosymbiotic: it descends from a once
free-living bacterium captured by an ancestral eukaryote roughly two
billion years ago, an account first defended in detail by Lynn Sagan
(later Margulis) in 1967 \cite{paper:v6c01-margulis-1967} and now
broadly confirmed by sequence comparison.

\begin{table}[htbp]
\centering
\caption[Canonical organelle inventory]{Canonical inventory of
eukaryotic organelles with their characteristic size, copy number per
mammalian cell, and primary function. Values are order-of-magnitude
defaults; particular cell types may differ by factors of ten or more.
Compiled from canonical reference works
\cite{text:v6c01-alberts-mbc-2022,text:v6c01-bionumbers-2015}.}
\label{tab:vol06:ch01:organelles}
\begin{tabular}{p{2.7cm} p{2.0cm} p{1.6cm} p{6.5cm}}
\hline
\textbf{Organelle} & \textbf{Size ($\mu$m)} &
\textbf{Copies/cell} & \textbf{Primary function} \\
\hline
Nucleus & $5$--$10$ & $1$ &
  Genome storage; transcription; mRNA export. \\
Mitochondrion & $0.5$--$1.5$ & $10^2$--$10^4$ &
  ATP synthesis by oxidative phosphorylation; calcium buffering;
  apoptosis initiation. \\
Endoplasmic reticulum & network & $1$ &
  Membrane-protein and lipid synthesis; quality control;
  calcium store. \\
Golgi apparatus & $3$--$10$ stacks & $10$--$100$ &
  Post-translational modification; protein sorting; vesicle
  generation. \\
Lysosome & $0.3$--$0.6$ & $\sim 300$ &
  Degradation of internalised and autophagic material; pH $\approx 4.5$. \\
Peroxisome & $0.1$--$0.5$ & $\sim 500$ &
  Long-chain fatty-acid oxidation; H$_2$O$_2$ metabolism. \\
Chloroplast & $5$--$10$ & $10$--$100$ &
  Photosynthesis; CO$_2$ fixation; starch storage (plants only). \\
Vacuole & up to cell-sized & $1$ &
  Turgor; storage; degradation (plants and fungi only). \\
\hline
\end{tabular}
\end{table}

\engterm{The cytoskeleton}. The interior of a eukaryotic cell is
mechanically organised by three classes of protein filament. The
\engterm{actin filaments} (microfilaments) are $7\,\si{\nano\meter}$
in diameter, dynamic on second-to-minute timescales, and form the
cortical network that gives the cell its shape, the contractile ring
that pinches the cell in two during cytokinesis, and the filaments
that drive cell crawling. The \engterm{intermediate filaments}, $10$
to $12\,\si{\nano\meter}$ in diameter, are mechanically the stiffest
of the three; they form the nuclear envelope's structural lamina, the
keratin network of skin cells, and the neurofilaments of neurons.
The \engterm{microtubules} are $25\,\si{\nano\meter}$ hollow tubes
built from tubulin dimers; they form the mitotic spindle that
segregates chromosomes, the tracks along which organelles and vesicles
are transported by the motor proteins kinesin and dynein (Chapter~\ref{vol06:ch04}),
and the structural core of cilia and flagella.
\Cref{fig:vol06:ch01:cytoskeleton} compares the three classes.

% Figure: three cytoskeletal filaments (actin, microtubules,
% intermediate filaments) drawn schematically with their characteristic
% diameter, structural unit, and rough mechanical role labelled.
\begin{figure}[htbp]
\centering
\begin{tikzpicture}[
  fil/.style={very thick, draw=engblack},
  monomer/.style={circle, draw=engblue!80!black, fill=engblue!30,
    inner sep=1.4pt, minimum size=0.34cm},
  tubmer/.style={circle, draw=engred!80!black, fill=engred!25,
    inner sep=1.2pt, minimum size=0.34cm},
  intf/.style={rectangle, draw=engamber!80!black, fill=engamber!25,
    minimum width=0.20cm, minimum height=0.45cm},
  lbl/.style={font=\small, align=center},
  callout/.style={font=\scriptsize, align=left},
]

% Actin filament (top): two-stranded helix; 7 nm diameter
\node[lbl, anchor=south] at (-2.5, 1.6) {\textbf{Actin} \\ \scriptsize 7 nm};
\foreach \i in {0,1,...,12} {
  \pgfmathsetmacro{\x}{0.45*\i}
  \pgfmathsetmacro{\yu}{1.7 + 0.20*sin(\i*40)}
  \pgfmathsetmacro{\yl}{1.5 + 0.20*sin(\i*40 + 180)}
  \node[monomer] at (\x, \yu) {};
  \node[monomer] at (\x, \yl) {};
}
\node[callout, anchor=west] at (6.4, 1.6)
  {dynamic: seconds to minutes \\ cortical network; cytokinesis; \\
   crawling};

% Microtubule (middle): hollow tube; 25 nm diameter
\node[lbl, anchor=south] at (-2.5, -0.3) {\textbf{Microtubule} \\ \scriptsize 25 nm};
\draw[fil, fill=engred!10, rounded corners=2pt]
  (0, -0.55) rectangle (5.4, 0.15);
\foreach \i in {0,1,...,11} {
  \pgfmathsetmacro{\x}{0.20 + 0.45*\i}
  \node[tubmer] at (\x, 0.05) {};
  \node[tubmer] at (\x, -0.45) {};
}
\node[callout, anchor=west] at (6.4, -0.2)
  {polarised; rigid \\ tracks for kinesin / dynein; \\ mitotic spindle};

% Intermediate filament (bottom): coiled-coil rope; 10 nm diameter
\node[lbl, anchor=south] at (-2.5, -2.0) {\textbf{Intermediate} \\ \scriptsize 10 nm};
\foreach \i in {0,1,...,18} {
  \pgfmathsetmacro{\x}{0.30*\i}
  \node[intf] at (\x, -1.7) {};
}
\node[callout, anchor=west] at (6.4, -1.7)
  {stiffest; non-polar \\ keratin; lamins; \\ neurofilaments};

% Diameter scale bar
\draw[|<->|, thick] (5.7, -2.5) -- (5.95, -2.5)
  node[midway, below, font=\scriptsize] {};

\end{tikzpicture}
\caption[Three cytoskeletal filament classes]{The three cytoskeletal
filament classes of eukaryotic cells, drawn schematically at relative
diameter (not to scale). \textbf{Actin} (7 nm) is dynamic on
second-to-minute timescales and forms the cortical network, the
contractile ring of cytokinesis, and the filaments of cell crawling.
\textbf{Microtubules} (25 nm hollow tubes) are stiff and polarised;
they form the mitotic spindle and the tracks along which kinesin and
dynein transport organelles and vesicles. \textbf{Intermediate
filaments} (10 nm) are the stiffest and the only non-polar class;
they include keratins, neurofilaments, and the nuclear-envelope
lamins. The three classes together set the cell's mechanical
properties \cite[ch.~16]{text:v6c01-alberts-mbc-2022}.}
\label{fig:vol06:ch01:cytoskeleton}
\end{figure}


The cytoskeleton sets the cell's mechanical properties. The cortical
actin network gives an animal cell a tension of approximately
$10^{-3}\,\text{N}/\text{m}$ \cite{paper:v6c01-tinevez-2009} and a
bulk elastic modulus of $100$ to $1000\,\text{Pa}$
\cite{text:v6c01-howard-mechanics-2001}. The combined cytoskeletal
network behaves on long timescales like a viscoelastic gel; on short
timescales, like an elastic solid. Persistence lengths characterise
the stiffness of each filament class: actin has a persistence length
of $\sim 15\,\si{\micro\meter}$, intermediate filaments
$\sim 1\,\si{\micro\meter}$, and microtubules
$\sim 5\,\text{mm}$ (microtubules behave as effectively rigid rods on
the scale of a cell). Chapter~\ref{vol06:ch06} will return to cellular
mechanics with the constitutive models the gel metaphor compresses.

\begin{archetype}[Interface and control at the cell scale]
The cell exhibits the interface and control archetypes in their
canonical biological form. The plasma membrane is an
\engterm{engineered interface} that separates two chemically distinct
compartments, supports selective transport across the boundary,
maintains the disequilibrium that defines life, and self-heals on
small disruption. The internal network of organelles, signalling
molecules, and gene-regulatory loops constitutes a
\engterm{distributed control system} with sensing (receptors and
metabolic sensors), actuation (transcription factors and post-translational
modifications), and feedback (homeostatic loops that maintain
intracellular pH, ion concentrations, energy charge, and cell volume
against perturbations). The two archetypes will recur in every
chapter of Volume VI: the interface, scaled up, becomes tissue and
organ design (Chapter~\ref{vol06:ch05}); the control system, scaled
up, becomes nervous and immune systems (Chapter~\ref{vol06:ch07}).
The atomic unit of both is the cell.
\end{archetype}

\section{Prokaryotic vs eukaryotic}\pathtag{standard}

Cells divide into two structural classes. \engterm{Prokaryotic} cells
lack a membrane-bounded nucleus; their DNA sits in the cytoplasm in a
loosely organised region called the nucleoid. The prokaryotes are the
bacteria and the archaea; their typical size is $1$ to
$5\,\si{\micro\meter}$. \engterm{Eukaryotic} cells have a
membrane-bounded nucleus and a full set of membrane-bounded organelles;
they range from $10$ to $100\,\si{\micro\meter}$ for most cell types,
and include all animals, plants, fungi, and protists. The transition
between the two classes is the largest single structural innovation in
the history of life, and it occurred approximately $2 \times 10^9$
years ago, an order of magnitude more recently than the origin of life
itself. \Cref{fig:vol06:ch01:prokaryote-eukaryote} draws the two
side by side.

% Figure: side-by-side schematic of a prokaryotic and a eukaryotic
% cell, with the canonical organelles labelled. Drawn approximately to
% relative scale on a log axis (linear-scale comparison would render
% the prokaryote invisibly small).
\begin{figure}[htbp]
\centering
\begin{tikzpicture}[
  cellwall/.style={very thick, draw=engblack, fill=engblue!8},
  membrane/.style={thick, draw=engred, fill=engblue!12},
  organelle/.style={thick, draw=engamber!90!black, fill=engamber!20},
  nuc/.style={thick, draw=engblue!80!black, fill=engblue!25},
  lbl/.style={font=\scriptsize, anchor=west},
  ttl/.style={font=\small\bfseries, align=center},
]

% Prokaryotic cell (left)
\node[ttl] at (1.6, 4.3) {Prokaryote \\ \scriptsize $\sim 1$--$5\,\mu$m};

% Outer wall (peptidoglycan)
\draw[cellwall] (1.6, 1.6) ellipse (1.5 and 1.0);
% Inner plasma membrane
\draw[membrane] (1.6, 1.6) ellipse (1.35 and 0.85);
% Nucleoid (no envelope)
\draw[densely dashed, very thick, draw=engblue!80!black]
  (1.6, 1.6) ellipse (0.55 and 0.35);
% Ribosomes
\foreach \p in {(0.5,1.4), (0.7,1.9), (2.6,1.8), (2.4,1.4),
                (1.8,1.1), (1.2,2.1), (2.2,2.0)} {
  \fill[engred] \p circle (0.05);
}
% Flagellum
\draw[engblack, thick, decorate, decoration={snake, amplitude=2pt,
  segment length=6pt}] (3.1, 1.5) -- (4.3, 1.5);

% Prokaryote labels (callouts)
\node[lbl] at (-0.8, 0.4) {peptidoglycan wall};
\draw[thin, draw=engblack] (0.5, 0.6) -- (0.7, 1.0);
\node[lbl] at (3.4, 2.6) {nucleoid (no envelope)};
\draw[thin, draw=engblack] (3.4, 2.55) -- (2.0, 1.85);
\node[lbl] at (3.4, 0.7) {ribosomes};
\draw[thin, draw=engblack] (3.4, 0.7) -- (2.4, 1.35);
\node[lbl] at (4.2, 1.5) {flagellum};

% Eukaryotic cell (right)
\node[ttl] at (10.5, 4.3) {Eukaryote \\ \scriptsize $\sim 10$--$100\,\mu$m};
% Plasma membrane (no wall in animal cells)
\draw[membrane] (10.5, 1.6) ellipse (3.0 and 2.2);
% Nucleus with envelope
\draw[nuc] (9.8, 2.2) ellipse (1.0 and 0.7);
% Nucleolus
\fill[engblue!60!black] (9.5, 2.2) ellipse (0.25 and 0.18);
% Mitochondria (3 instances)
\draw[organelle] (11.3, 2.6) ellipse (0.45 and 0.20);
\draw[engamber!90!black, thick] (10.9, 2.55) -- (11.7, 2.65);
\draw[organelle] (8.7, 1.2) ellipse (0.45 and 0.20);
\draw[engamber!90!black, thick] (8.3, 1.15) -- (9.1, 1.25);
\draw[organelle] (11.8, 1.0) ellipse (0.45 and 0.20);
\draw[engamber!90!black, thick] (11.4, 0.95) -- (12.2, 1.05);
% Endoplasmic reticulum (curved tubes)
\draw[thick, draw=engred!80!black] (10.2, 0.9)
  to[out=20, in=200] (11.5, 1.4)
  to[out=20, in=200] (12.2, 1.5);
\draw[thick, draw=engred!80!black] (10.0, 0.6)
  to[out=10, in=180] (11.3, 0.9);
% Golgi (stack of curved discs)
\draw[thick, draw=engamber!80!black] (9.5, 0.5) to[bend right=20] (10.4, 0.5);
\draw[thick, draw=engamber!80!black] (9.5, 0.35) to[bend right=20] (10.4, 0.35);
\draw[thick, draw=engamber!80!black] (9.5, 0.20) to[bend right=20] (10.4, 0.20);
% Cytoskeleton (sketchy filaments)
\draw[thin, draw=enggray] (8.3, 3.0) -- (12.3, 1.7);
\draw[thin, draw=enggray] (8.5, 2.0) -- (12.5, 2.6);
\draw[thin, draw=enggray] (9.5, 3.2) -- (11.8, 0.7);

% Eukaryote labels
\node[lbl, anchor=east] at (8.0, 0.4) {plasma membrane};
\draw[thin, draw=engblack] (8.0, 0.45) -- (8.4, 0.6);
\node[lbl, anchor=east] at (8.0, 2.2) {nucleus};
\draw[thin, draw=engblack] (8.0, 2.2) -- (8.85, 2.2);
\node[lbl, anchor=east] at (8.0, 2.7) {nucleolus};
\draw[thin, draw=engblack] (8.0, 2.65) -- (9.3, 2.2);
\node[lbl, anchor=west] at (12.0, 3.0) {mitochondrion};
\draw[thin, draw=engblack] (12.0, 3.0) -- (11.6, 2.7);
\node[lbl, anchor=west] at (12.6, 1.5) {ER};
\draw[thin, draw=engblack] (12.6, 1.5) -- (12.3, 1.5);
\node[lbl, anchor=west] at (10.6, -0.05) {Golgi};
\draw[thin, draw=engblack] (10.6, -0.05) -- (10.2, 0.2);
\node[lbl, anchor=west] at (12.6, 2.4) {cytoskeleton};
\draw[thin, draw=engblack] (12.6, 2.4) -- (12.3, 2.4);

\end{tikzpicture}
\caption[Prokaryote and eukaryote, side by side]{A prokaryotic cell
(left, drawn approximately $1$--$5\,\si{\micro\meter}$ across) and a
eukaryotic cell (right, drawn approximately $10$--$100\,\si{\micro\meter}$
across). The prokaryote has a peptidoglycan cell wall, a single plasma
membrane, ribosomes free in the cytoplasm, a single circular
chromosome in a nucleoid region without a membrane envelope, and
appendages such as flagella. The eukaryote has a membrane-bounded
nucleus with a nucleolus, mitochondria with their double membrane,
endoplasmic reticulum, Golgi apparatus, and a cytoskeleton that
organises the interior \cite[ch.~1]{text:v6c01-alberts-mbc-2022}.
Scale within each panel is approximately consistent; the two panels
are at incommensurate scales (linear comparison would render the
prokaryote invisible).}
\label{fig:vol06:ch01:prokaryote-eukaryote}
\end{figure}


\engterm{Size-scaling differences}. The size difference between
prokaryote and eukaryote (a factor of $10$ in linear dimension, a
factor of $1000$ in volume) corresponds to a difference in surface-to-volume
ratio of a factor of $10$. The prokaryotic cell has enough
boundary area, relative to its volume, that all metabolic transport
can occur across the plasma membrane; the eukaryotic cell does not.
The eukaryotic cell solves the deficit by internalising membrane:
mitochondria, ER, Golgi, lysosomes, and vesicles together expose
membrane surface area of $20$ to $30$ times the plasma membrane area.
The internal membrane is also chemically specialised, each compartment
maintaining its own pH, ionic composition, and protein complement.

\engterm{Genome size and information capacity}. A typical bacterial
genome is $10^6$ to $10^7$ base pairs and encodes $1000$ to $10\,000$
proteins. A typical eukaryotic genome ranges from $10^7$ (yeast) to
$10^{11}$ base pairs (some plants, salamanders, and lungfish), with
the human genome at $3.2 \times 10^9$ base pairs encoding roughly
$20\,000$ protein-coding genes. The mismatch between bacterial
parsimony (high coding density, little non-coding sequence) and
eukaryotic profligacy (most of the genome is non-coding,
regulatory, repetitive, or transposable) is one of the central
puzzles of comparative genomics and will return in Chapter~\ref{vol06:ch03}.

\engterm{Generation time}. A bacterium can divide in $20\,\text{minutes}$
under optimal laboratory conditions; \emph{Escherichia coli} doubles
every $30\,\text{minutes}$ on rich medium. A mammalian cell in
culture doubles in $16$ to $24\,\text{hours}$. A plant cell may take
several days. The generation-time difference sets the timescales on
which prokaryotic and eukaryotic engineering operate:
biotechnological yields from \emph{E.~coli} arrive in hours, from
\emph{Saccharomyces cerevisiae} (yeast) in a day, from mammalian
cell culture in a week. Chapter~\ref{vol06:ch10} on biomanufacturing
will use the generation-time ranking as a primary cost driver.

\engterm{The nuclear envelope as gatekeeper}. The eukaryote's
double-membrane nuclear envelope is perforated by nuclear pore
complexes (NPCs) that select what crosses into and out of the
nucleus. Small molecules (Stokes radius $\lesssim 4\,\text{nm}$) pass
freely; larger cargoes carry a nuclear-localisation or nuclear-export
signal recognised by importin or exportin proteins, which use the
RanGTP gradient as an energy source to translocate the cargo
\cite{paper:v6c01-paine-1975}. The arrangement is a textbook case of
sieve-plus-active-transport at the boundary, with the boundary
selectivity coupled to a primary metabolic gradient. Prokaryotes
have no analogue: their genome and ribosomes share a single
compartment, so transcription and translation are continuous in
space and in time.

\engterm{Engineering tractability}. Prokaryotic cells are simpler,
faster, and cheaper to engineer; eukaryotic cells provide
post-translational modifications (glycosylation, complex disulfide
bonding, folding chaperones) that some product molecules require.
Insulin was first produced biotechnologically in bacterial
\emph{E.~coli} in 1978; many therapeutic antibodies are produced in
mammalian Chinese Hamster Ovary (CHO) cells because the antibodies
require the eukaryotic glycosylation machinery. The choice between
prokaryote and eukaryote for a given biomanufacturing target is set
by the molecular complexity of the target, with the prokaryote chosen
when feasible and the eukaryote chosen when necessary.

\section{The cell as control loop}\pathtag{core}

A cell is a control system. The reading is not metaphorical: a cell
contains sensors that report the state of its environment, actuators
that change its internal or external state, and feedback loops that
regulate the actuators based on the sensor outputs. The vocabulary of
classical control theory (set point, error signal, integral action,
gain, time constant, stability margin) maps onto cellular biology
directly. Uri Alon's \emph{Systems Biology} text is the canonical
working reference for the engineering of biological control loops;
the rest of this section follows his framing
\cite[ch.~2--5]{text:alon-systems-biology}, with Dennis Bray's earlier
framing of intracellular signalling as biological computation
\cite{paper:v6c01-bray-1995} as the historical pivot point.
\Cref{fig:vol06:ch01:control-loop} sketches the loop in
block-diagram form.

% Figure: the cell as a control loop. Sensor -> Controller -> Actuator
% -> Output -> Feedback. The canonical block diagram of classical
% control engineering, redrawn with cellular components named at each
% block.
\begin{figure}[htbp]
\centering
\begin{tikzpicture}[
  block/.style={draw=engblack, thick, rectangle, minimum width=2.6cm,
    minimum height=1.1cm, align=center, font=\small, fill=white},
  sense/.style={block, fill=engblue!10},
  ctrl/.style={block, fill=engamber!15},
  act/.style={block, fill=engred!10},
  sum/.style={draw=engblack, thick, circle, minimum size=0.6cm,
    inner sep=0pt, font=\small},
  arr/.style={-{Latex[length=2.2mm]}, thick, draw=engblack},
  lbl/.style={font=\scriptsize\itshape, align=center},
]
% Top row: forward path
\node[sum] (sum) at (0,0) {$\Sigma$};
\node[ctrl] (ctrl) at (3,0) {Controller \\ \scriptsize (transcription \\ \scriptsize factor cascade)};
\node[act] (act) at (6.7,0) {Actuator \\ \scriptsize (enzyme; channel; \\ \scriptsize pump)};
\node[block, fill=englime!15] (plant) at (10.3,0) {Process \\ \scriptsize (the cell's \\ \scriptsize chemistry)};

% Output
\node[block, fill=englime!20] (out) at (10.3,-2.4) {Output \\ \scriptsize (concentration; \\ \scriptsize voltage; rate)};

% Sensor
\node[sense] (sens) at (3,-2.4) {Sensor \\ \scriptsize (receptor; \\ \scriptsize metabolic sensor)};

% Setpoint and disturbance
\node[lbl, anchor=east] (sp) at (-1.4,0) {setpoint \\ (genetic / \\ homeostatic)};
\node[lbl, anchor=south] (dist) at (10.3,1.4) {disturbance \\ (stress; toxin; \\ ligand)};

% Arrows forward
\draw[arr] (sp) -- (sum);
\draw[arr] (sum) -- (ctrl) node[midway, above, font=\scriptsize] {error};
\draw[arr] (ctrl) -- (act);
\draw[arr] (act) -- (plant);
\draw[arr] (dist) -- (plant);
\draw[arr] (plant) -- (out) node[midway, right, font=\scriptsize] {measured};

% Feedback path
\draw[arr] (out) -- (sens);
\draw[arr] (sens) -| node[near end, left, font=\scriptsize] {$-$} (sum.south);

\end{tikzpicture}
\caption[The cell as a control loop]{The cell as a control loop. A
setpoint (encoded genetically or homeostatically) is compared against
a sensor reading; the error drives a controller (a transcription
factor cascade, a second-messenger cascade, or a covalent modification)
that drives an actuator (an enzyme, an ion channel, a pump) that acts
on the process (the cell's internal chemistry). The output is measured
by a sensor whose reading closes the loop. Each block has a named
cellular implementation, and the cell's behaviour at the loop level is
governed by the same equations as the engineering analogue.}
\label{fig:vol06:ch01:control-loop}
\end{figure}


\engterm{Negative feedback}. The simplest cellular control loop is a
negative feedback regulating the concentration of some target
molecule against perturbation. A bacterial cell synthesising the amino
acid tryptophan illustrates the canonical case. The bacterial
\emph{trp} operon contains five genes that encode the enzymes for
tryptophan biosynthesis. The operon's expression is repressed by a
transcription factor that becomes active only when bound to
tryptophan. When intracellular tryptophan is high, the repressor binds
DNA and the operon is off; when tryptophan is low, the repressor is
inactive, the operon is on, more tryptophan is synthesised, and the
level rises back toward the set point. The loop is end-product
inhibition by transcriptional repression. The gain is set by the
binding constants of the repressor-DNA and repressor-tryptophan
interactions; the response time is set by the lifetime of the
repressor protein and the operon's mRNA.

\engterm{Feed-forward loops}. The bacterial transcription network
contains motifs that recur far more often than random graphs would
predict. The most common is the feed-forward loop: a regulator $X$
activates two targets $Y$ and $Z$, and $Y$ also activates $Z$. The
loop comes in eight variants depending on the signs of the three
edges. The coherent type-1 feed-forward loop (all three edges
activating) responds to a step-up of $X$ with a delayed activation of
$Z$ (because $Z$ requires both $X$ and $Y$, and $Y$ takes time to
accumulate) but to a step-down of $X$ with prompt deactivation of
$Z$. The asymmetric response acts as a noise filter: brief excursions
of $X$ do not propagate to $Z$, while sustained excursions do
\cite[ch.~4]{text:alon-systems-biology,paper:v6c01-shoval-alon-2010}.
The feed-forward loop is present in roughly half of the inducible
operons in \emph{E.~coli} and in many developmental gene-regulatory
networks; it is one of the volume's central examples of an
engineering motif read out of cellular biology.
\Cref{fig:vol06:ch01:feedforward-loop} draws the loop and its
asymmetric step response.

% Figure: coherent type-1 feed-forward loop. X activates Y; X and Y
% both activate Z; the loop filters brief excursions of X.
\begin{figure}[htbp]
\centering
\begin{tikzpicture}[
  node distance=2.6cm,
  reg/.style={circle, draw=engblack, very thick, fill=engblue!15,
    minimum size=1.0cm, font=\bfseries},
  act/.style={-{Latex[length=2.2mm]}, thick, draw=engblack},
  lbl/.style={font=\scriptsize, align=left},
]

% Topology (left)
\begin{scope}[shift={(0,0)}]
\node[reg] (x) at (0,2.6) {X};
\node[reg] (y) at (-1.4, 0.6) {Y};
\node[reg] (z) at (1.4, 0.6) {Z};
\draw[act] (x) -- (y);
\draw[act] (x) -- (z);
\draw[act] (y) -- (z);
\node[font=\small\bfseries] at (0, -0.6) {Topology};
\end{scope}

% Step-up response (middle)
\begin{scope}[shift={(5.5,0)}]
\draw[->, thick] (-0.2, 0) -- (3.6, 0);
\draw[->, thick] (-0.2, 0) -- (-0.2, 3.0);
% X step-up
\draw[engblue, very thick] (-0.2, 0.4) -- (0.6, 0.4)
  -- (0.6, 2.5) -- (3.4, 2.5);
\node[font=\scriptsize, engblue, anchor=west] at (3.5, 2.5) {X};
% Z delayed
\draw[engred, very thick]
  (-0.2, 0.4) -- (0.6, 0.4)
  to[out=15, in=180] (1.5, 0.5)
  to[out=0, in=200] (2.5, 2.0)
  to[out=20, in=180] (3.4, 2.3);
\node[font=\scriptsize, engred, anchor=west] at (3.5, 2.3) {Z};
\node[font=\scriptsize, anchor=south] at (1.5, 2.7) {delay};
\node[font=\small\bfseries] at (1.5, -0.6) {Step-up};
\node[font=\scriptsize, anchor=west] at (-0.2, 3.2) {amplitude};
\node[font=\scriptsize, anchor=west] at (3.6, -0.05) {time};
\end{scope}

% Step-down response (right)
\begin{scope}[shift={(10.5,0)}]
\draw[->, thick] (-0.2, 0) -- (3.6, 0);
\draw[->, thick] (-0.2, 0) -- (-0.2, 3.0);
% X step-down
\draw[engblue, very thick] (-0.2, 2.5) -- (1.0, 2.5)
  -- (1.0, 0.4) -- (3.4, 0.4);
\node[font=\scriptsize, engblue, anchor=west] at (3.5, 0.4) {X};
% Z drops promptly
\draw[engred, very thick]
  (-0.2, 2.3) -- (1.0, 2.3)
  to[out=-30, in=180] (1.4, 0.5) -- (3.4, 0.5);
\node[font=\scriptsize, engred, anchor=west] at (3.5, 0.5) {Z};
\node[font=\scriptsize, anchor=south] at (1.4, 1.5) {prompt};
\node[font=\small\bfseries] at (1.5, -0.6) {Step-down};
\node[font=\scriptsize, anchor=west] at (3.6, -0.05) {time};
\end{scope}

\end{tikzpicture}
\caption[Coherent type-1 feed-forward loop]{The coherent type-1
feed-forward loop is the most common motif in the \emph{E.\ coli}
transcription network \cite{paper:v6c01-shoval-alon-2010}. The
regulator $X$ activates both $Y$ and the output $Z$, and $Y$ in turn
activates $Z$ (left). The loop responds asymmetrically: a step-up of
$X$ produces a \emph{delayed} activation of $Z$ (middle), because
$Z$ requires both $X$ and $Y$ and $Y$ takes time to accumulate; a
step-down of $X$ produces a \emph{prompt} deactivation of $Z$ (right).
The asymmetry implements a sign-sensitive delay that filters brief
excursions of $X$ without filtering brief drops.}
\label{fig:vol06:ch01:feedforward-loop}
\end{figure}


\engterm{Bistability and switches}. A cell needs to choose between
discrete states (replicating versus quiescent, undifferentiated versus
committed, healthy versus apoptotic) and to maintain the chosen state
in the face of noise. The engineering device for the choice is a
\engterm{bistable switch}: a regulatory network with two stable
steady states separated by an unstable intermediate. The simplest
bistable circuit is mutual repression: two transcription factors $X$
and $Y$ each repress the other. If $X$ wins, $Y$ is low and the
repression of $X$ is weak; the state with high $X$, low $Y$ is stable.
Symmetrically, the state with low $X$, high $Y$ is stable. A
sufficiently strong perturbation flips the system; small perturbations
do not. The phage lambda lysis/lysogeny switch is the textbook example
\cite[ch.~5]{text:alon-systems-biology}. Cell-fate commitment in
mammalian development uses analogous circuits, and the apoptotic
decision (section~1.6) is bistable in the same engineering sense.

\engterm{Stochastic noise}. Cellular regulation runs on small numbers
of molecules: a typical bacterial transcription factor is present at
$10$ to $100$ molecules per cell, and an mRNA at $1$ to $10$ copies.
The stochastic biochemistry of small numbers produces appreciable
cell-to-cell variation that the cell's regulatory architecture must
either suppress or exploit. Elowitz and colleagues' 2002 two-colour
fluorescent-reporter experiment in \emph{E.~coli} cleanly separated
\engterm{intrinsic noise} (the stochastic biochemistry within one
cell) from \engterm{extrinsic noise} (cell-to-cell variation in
regulatory factors), establishing a quantitative framework for the
noise budget of a gene-regulatory circuit
\cite{paper:v6c01-elowitz-2002}. The intrinsic-noise floor scales
as $1/\sqrt{N}$ with the average copy number $N$; the extrinsic-noise
contribution does not, and is the dominant term for highly expressed
genes.

\engterm{The cell cycle as a control system}. The eukaryotic cell
cycle is a sequence of four phases (G1, S, G2, M) regulated by
cyclin-dependent kinases. The kinases activate downstream effectors
that drive DNA replication and mitosis; the cyclins are synthesised
and degraded in patterns that produce phase-specific kinase activity.
The cycle is gated at three checkpoints (G1/S, G2/M, and the spindle
assembly checkpoint at M) that halt the cycle if DNA is damaged or
chromosomes are misaligned. The architecture, developed by Hartwell,
Hunt, and Nurse (Nobel Prize 2001), is a relaxation oscillator with
state-dependent inputs \cite{paper:v6c01-nurse-2001}. The engineering
of the cell cycle is the subject of Chapter~\ref{vol06:ch03}; the
failure mode (loss of checkpoint integrity, leading to genome
instability and cancer) is the subject of that chapter's failure
section.

\engterm{Adaptation}. A bacterium swimming up a chemical gradient
needs to discriminate between an increase in attractant concentration
(``swim further'') and a high but constant concentration (``the
gradient is gone, change direction''). The discrimination is
\engterm{exact adaptation}: the steady-state response to a step input
returns to a fixed baseline regardless of the input's magnitude. The
bacterial chemotaxis system achieves exact adaptation by an
integral-feedback mechanism whose structure is identical to the
integral-controller pole of a PID controller
\cite{paper:v6c01-yi-2000,text:alon-systems-biology}. The reading is
exact: the bacterium's adaptation circuit is, in control-theoretic
terms, an integral controller, with the integration performed by the
slow methylation of chemoreceptor proteins. The amplification gain
between single-receptor occupancy and flagellar reversal can exceed
a factor of $50$ \cite{paper:v6c01-bray-2002}; the gain-times-bandwidth
product of the chemotaxis controller is comparable to that of
engineered analog controllers of similar response time.

\begin{principle}[Cellular regulation is control engineering]
The mathematical structures of cellular regulation are the mathematical
structures of control engineering. Negative feedback, feed-forward
loops, bistable switches, oscillators, and integral controllers all
appear in named cellular circuits with the same governing equations
as their engineering analogues. The cell does not happen to use the
same mathematics; the cell is doing control engineering, and the
engineer reading the cell is reading control engineering performed by
evolution. Volume IX develops control theory in its general form;
the present chapter establishes its biological instantiation.
\end{principle}

\section{Worked examples}\pathtag{standard}

The cell types of biology are diverse, but for our engineering
purposes the variety is tractable through four worked examples that
exhibit, between them, the principal cellular archetypes. Each is
treated as an engineered object: what it does, what it costs, where
it fails. \Cref{fig:vol06:ch01:worked-examples} shows the four
side by side at relative shape.

% Figure: four-panel comparison of the chapter's worked-example cells:
% red blood cell, neuron, plant cell, bacterium. Each panel shows the
% cell's characteristic shape and one or two engineered features.
\begin{figure}[htbp]
\centering
\begin{tikzpicture}[
  membrane/.style={very thick, draw=engred, fill=engblue!10},
  organelle/.style={thick, draw=engamber!90!black, fill=engamber!20},
  nuc/.style={thick, draw=engblue!80!black, fill=engblue!25},
  cellwall/.style={very thick, draw=engblack, fill=engblue!8},
  ttl/.style={font=\small\bfseries, align=center},
  lbl/.style={font=\scriptsize, align=center},
]

% Panel 1: Red blood cell (biconcave disc)
\begin{scope}[shift={(0,0)}]
\node[ttl] at (1.5, 3.0) {RBC \\ \scriptsize $7.5\,\mu$m, biconcave};
\draw[membrane] (0.0, 1.6)
  to[out=90, in=180] (1.5, 2.4)
  to[out=0, in=90] (3.0, 1.6)
  to[out=-90, in=0] (1.5, 0.7)
  to[out=180, in=-90] (0.0, 1.6);
\draw[membrane, fill=engblue!25] (0.6, 1.6) ellipse (0.4 and 0.18);
\draw[membrane, fill=engblue!25] (2.4, 1.6) ellipse (0.4 and 0.18);
\node[lbl] at (1.5, 0.2) {no nucleus, no mitochondria \\ haemoglobin to 330 g/L};
\end{scope}

% Panel 2: Neuron
\begin{scope}[shift={(5.0,0)}]
\node[ttl] at (1.5, 3.0) {Neuron \\ \scriptsize soma $10$--$100\,\mu$m; axon up to $1$ m};
% Dendrites
\draw[membrane] (0.5, 2.2) -- (1.1, 1.85);
\draw[membrane] (0.3, 1.8) -- (1.1, 1.7);
\draw[membrane] (0.4, 1.4) -- (1.1, 1.5);
\draw[membrane] (0.7, 1.0) -- (1.2, 1.4);
% Soma
\draw[membrane] (1.5, 1.7) ellipse (0.5 and 0.45);
\draw[nuc] (1.5, 1.7) ellipse (0.22 and 0.18);
% Axon (long, with myelin segments)
\draw[membrane] (2.0, 1.7) -- (3.4, 1.7);
\foreach \x in {2.1, 2.5, 2.9} {
  \draw[organelle, fill=engamber!25]
    (\x, 1.6) rectangle (\x+0.3, 1.8);
}
% Terminal
\draw[membrane] (3.4, 1.7) -- (3.6, 1.85);
\draw[membrane] (3.4, 1.7) -- (3.6, 1.55);
\node[lbl] at (1.5, 0.4) {action potential 100 mV $\cdot$ 1 ms \\
  myelin reduces $C_m$ by 50--100$\times$};
\end{scope}

% Panel 3: Plant cell
\begin{scope}[shift={(0,-4.0)}]
\node[ttl] at (1.5, 3.0) {Plant cell \\ \scriptsize $10$--$100\,\mu$m, walled};
% Cell wall (outer thick boundary)
\draw[cellwall, line width=2.5pt] (0.2, 0.5) rectangle (2.8, 2.6);
% Plasma membrane just inside
\draw[membrane] (0.35, 0.65) rectangle (2.65, 2.45);
% Central vacuole (big, central)
\draw[thick, draw=engblue!70!black, fill=engblue!15]
  (1.5, 1.55) ellipse (0.85 and 0.65);
\node[lbl] at (1.5, 1.55) {vacuole};
% Chloroplasts
\draw[organelle, fill=englime!30, draw=englime!60!black]
  (0.6, 2.0) ellipse (0.20 and 0.12);
\draw[organelle, fill=englime!30, draw=englime!60!black]
  (2.4, 1.0) ellipse (0.20 and 0.12);
\draw[organelle, fill=englime!30, draw=englime!60!black]
  (0.7, 0.85) ellipse (0.20 and 0.12);
% Nucleus (offset by vacuole)
\draw[nuc] (2.35, 2.15) ellipse (0.25 and 0.20);
\node[lbl] at (1.5, 0.0) {turgor 0.5--1.0 MPa; \\ wall as pressure vessel};
\end{scope}

% Panel 4: Bacterium (rod-shaped E. coli)
\begin{scope}[shift={(5.0,-4.0)}]
\node[ttl] at (1.5, 3.0) {Bacterium \\ \scriptsize \emph{E.\ coli}, $2 \times 0.5\,\mu$m};
% Cell envelope (rod shape)
\draw[cellwall, line width=1.5pt, rounded corners=8pt]
  (0.4, 1.4) rectangle (2.6, 2.0);
\draw[membrane, rounded corners=6pt]
  (0.55, 1.5) rectangle (2.45, 1.9);
% Nucleoid
\draw[densely dashed, very thick, draw=engblue!80!black]
  (1.5, 1.7) ellipse (0.6 and 0.18);
% Ribosomes
\foreach \p in {(0.8,1.6), (1.0,1.85), (2.1,1.6), (1.9,1.85),
                (0.7,1.78)} {
  \fill[engred] \p circle (0.04);
}
% Flagella (multiple)
\draw[engblack, thick, decorate, decoration={snake, amplitude=1.5pt,
  segment length=4pt}] (2.6, 1.7) -- (3.5, 1.6);
\draw[engblack, thick, decorate, decoration={snake, amplitude=1.5pt,
  segment length=4pt}] (2.6, 1.55) -- (3.5, 1.4);
\node[lbl] at (1.5, 0.4) {single chromosome 4.6 Mb \\ doubles in 30 min};
\end{scope}

\end{tikzpicture}
\caption[Four worked-example cells]{The four worked-example cells of
section 1.5, drawn at relative shape but not to relative scale. The
\textbf{red blood cell} (top left) is biconcave, anucleate, and packs
haemoglobin to one-third its mass. The \textbf{neuron} (top right) is
specialised for long-distance electrical signalling along its
myelinated axon. The \textbf{plant cell} (bottom left) is a
pressure-vessel design with a rigid cellulose wall containing a
turgid vacuole. The \textbf{bacterium} (bottom right) is a small rod
without internal membrane organelles, with a single circular
chromosome and rotary flagella.}
\label{fig:vol06:ch01:worked-examples}
\end{figure}


\engterm{Worked example 1: the red blood cell}. A mammalian
erythrocyte is the simplest cell in the body. It is a biconcave disc,
$7.5\,\si{\micro\meter}$ across by $2\,\si{\micro\meter}$ at the rim
and $1\,\si{\micro\meter}$ in the middle, with no nucleus and no
mitochondria. It contains haemoglobin at a concentration of
$330\,\text{g}/\text{L}$, more than a third of its mass. Its only job
is to carry oxygen from the lung to the tissues and to carry carbon
dioxide back.

The cell is engineered for the job by three structural choices. First,
the biconcave shape maximises the surface-to-volume ratio for a given
volume (the surface-to-volume ratio is $1.5$ times that of a sphere of
the same volume), which speeds diffusive oxygen exchange across the
plasma membrane. Second, the absence of a nucleus and organelles
maximises the haemoglobin concentration that can fit inside the cell.
Third, the cell's spectrin-actin cortical cytoskeleton makes the cell
mechanically pliable: it deforms to squeeze through capillaries
$3\,\si{\micro\meter}$ wide (narrower than the cell's resting
diameter) and returns to its biconcave shape downstream. The
mechanical flexibility is achieved at the cost of biosynthetic
capability: a mature erythrocyte cannot synthesise new proteins, repair
damaged components, or divide, and it lives for only $120\,\text{days}$
before being recycled by macrophages in the spleen and liver.

Numerical reading. The body contains $2.5 \times 10^{13}$
erythrocytes, a total of $7 \times 10^{17}$ haemoglobin tetramers,
each capable of binding four oxygen molecules. The lungs load
oxygen into haemoglobin at a partial pressure of $100\,\text{mmHg}$,
where saturation is approximately $98\,\%$; the tissues unload at
$40\,\text{mmHg}$, where saturation is approximately $75\,\%$. The
$23\,\%$ saturation drop, applied to $7 \times 10^{17}$ tetramers
$\times 4$ binding sites, releases $6.5 \times 10^{17}$ oxygen
molecules per litre of blood per pass, or about $200\,\text{mL}$ of
oxygen per litre of blood at standard conditions. The cardiac
output of $5\,\text{L}/\text{min}$ delivers $1\,\text{L}$ of oxygen
per minute, matching the body's resting consumption. The cell is
engineered for the throughput; nothing more.

\engterm{The historical RBC: Henrietta Lacks and HeLa}. The mention
of red blood cells brings into view a parallel question about
provenance and consent that runs through the volume. In 1951
Henrietta Lacks, a Black woman receiving treatment for cervical
cancer at the Johns Hopkins Hospital, had a tumour biopsy taken
without her knowledge or consent. The cells from that biopsy
established the first immortalised human cell line, ``HeLa,'' which
has been propagated continuously in laboratories worldwide for over
seven decades and used in tens of thousands of published studies. The
cells underlie the development of the polio vaccine, of many
chemotherapeutic agents, and of much of modern molecular cell biology.
Lacks died without ever being told her cells had been taken, and her
family was not informed of the cell line's existence or commercial
value until the 1970s. The HeLa story is the canonical case for the
non-technical engineering questions of consent, provenance, and
ownership of biological material that any tissue-derived therapy or
research programme must answer. We name it here without giving the
case its full treatment; the cell-line-provenance question returns
in the volume's biomanufacturing chapter
(Chapter~\ref{vol06:ch10}).

\engterm{Worked example 2: the neuron}. A vertebrate neuron is a
specialised cell for the long-distance propagation of electrical
signals. It has a cell body (soma) of $10$ to $100\,\si{\micro\meter}$,
a long output projection (axon) that may extend more than a metre in
large animals, and a branching input structure (dendritic tree) that
receives signals from up to thousands of other neurons. The axon
carries action potentials, voltage pulses of approximately
$100\,\text{mV}$ amplitude and $1\,\text{ms}$ duration, at velocities
that range from $1\,\text{m}/\text{s}$ in unmyelinated fibres to
$100\,\text{m}/\text{s}$ in the largest myelinated fibres. Synapses
at the axon's distal end couple the electrical signal to the next
neuron through chemical neurotransmitter release. The biophysical
description that quantitatively predicts the action-potential
waveform from independently measured channel kinetics is the 1952
Hodgkin-Huxley model \cite{paper:v6c01-hodgkin-huxley-1952}, one of
the canonical demonstrations of predictive engineering biology.

The cell is engineered for signal propagation. The axon is a
cylindrical conducting cable of about $0.2$ to $20\,\si{\micro\meter}$
diameter, surrounded by an insulating myelin sheath in vertebrate
nervous systems. The capacitance of the bare membrane,
$0.9\,\si{\micro\farad}/\text{cm}^2$, would charge slowly under the
ionic currents available; the myelin sheath, which is several wrapped
layers of plasma membrane from a supporting glial cell, reduces the
effective capacitance by a factor of $50$ to $100$ and increases the
membrane resistance similarly. The conduction velocity scales with
the square root of axon diameter in unmyelinated fibres and linearly
with diameter in myelinated fibres; the squared-versus-linear scaling
is the engineering case for myelin and is developed in detail in
Chapter~\ref{vol06:ch07}.

The metabolic cost of a single action potential, summed over the
sodium that enters and must be pumped back out, is roughly
$10^9$ ATP per action potential per centimetre of axon
\cite{paper:v6c01-attwell-laughlin-2001}. The human brain consumes
about $20\,\%$ of the body's resting metabolic rate while constituting
$2\,\%$ of its mass. The brain is the body's most energy-dense organ;
the neurons are why. The Na$^+$/K$^+$ ATPase, discovered by Skou in
1957 \cite{paper:v6c01-skou-1957}, is the molecular workhorse that
restores the ion gradients after each spike; the pump consumes roughly
half of the brain's resting ATP turnover
\cite{paper:v6c01-erecinska-1989}. \Cref{fig:vol06:ch01:na-k-pump}
sketches the pump's three-Na-out / two-K-in / one-ATP cycle.

% Figure: the sodium-potassium ATPase as a schematic, with the
% intracellular and extracellular ion concentrations labelled. Shows
% three Na+ out, two K+ in, one ATP hydrolysed.
\begin{figure}[htbp]
\centering
\begin{tikzpicture}[
  membrane/.style={very thick, draw=engred},
  pump/.style={fill=engamber!30, draw=engamber!80!black, very thick,
    rounded corners=2pt},
  na/.style={circle, fill=engblue!70, draw=engblue!90!black,
    inner sep=1.4pt, minimum size=0.36cm, font=\scriptsize,
    text=white},
  k/.style={circle, fill=engred!70, draw=engred!90!black,
    inner sep=1.4pt, minimum size=0.36cm, font=\scriptsize,
    text=white},
  arr/.style={-{Latex[length=2.0mm]}, thick, draw=engblack},
  lbl/.style={font=\scriptsize, align=left},
]

% Plasma membrane (two parallel lines)
\draw[membrane] (-1.0, 1.2) -- (10.0, 1.2);
\draw[membrane] (-1.0, 0.0) -- (10.0, 0.0);

% Pump in the middle
\draw[pump, rounded corners=4pt]
  (4.0, -0.3) -- (5.7, -0.3) -- (5.9, 1.5) -- (3.8, 1.5) -- cycle;
\node[font=\scriptsize\bfseries, align=center] at (4.85, 0.6) {Na$^+$/K$^+$ \\ ATPase};

% Ions on the intracellular side (Na+ inside in)
\node[na] at (3.0, 0.5) {Na};
\node[na] at (3.3, 0.2) {Na};
\node[na] at (2.8, 0.85) {Na};
% ATP
\node[font=\small\bfseries] at (3.1, -0.6) {ATP};

% Ions emerging extracellularly
\node[na] at (6.6, 1.85) {Na};
\node[na] at (7.0, 1.95) {Na};
\node[na] at (7.3, 1.65) {Na};

% K+ on the extracellular side (going in)
\node[k] at (7.0, 1.6) {K};
\node[k] at (6.6, 1.7) {K};

% K+ delivered intracellularly
\node[k] at (4.85, -0.8) {K};
\node[k] at (5.25, -0.65) {K};

% Arrows showing flux direction
\draw[arr] (3.4, 0.5) to[bend left=20]
  node[midway, above, font=\scriptsize] {$3 \times$ out} (5.5, 1.5);
\draw[arr] (4.8, 1.5) to[bend right=20]
  node[midway, below, font=\scriptsize] {$2 \times$ in} (4.4, 0.0);
\draw[arr] (3.4, -0.45) -- (3.95, -0.05)
  node[midway, below right, font=\scriptsize] {hydrolysed};

% Labels for compartments
\node[lbl, anchor=west] at (-0.5, 2.2) {\textbf{extracellular} \\
  $[\text{Na}^+]_o = 145$ mM \\ $[\text{K}^+]_o = 4$ mM};
\node[lbl, anchor=west] at (-0.5, -1.0) {\textbf{cytoplasm} \\
  $[\text{Na}^+]_i = 10$ mM \\ $[\text{K}^+]_i = 140$ mM};

\end{tikzpicture}
\caption[The sodium-potassium ATPase]{The sodium-potassium ATPase
(Na$^+$/K$^+$-ATPase) exchanges three intracellular sodium ions for
two extracellular potassium ions per molecule of ATP hydrolysed
\cite{paper:v6c01-skou-1957}. The net result is a current of one
positive charge out per cycle and the maintenance of the canonical
mammalian ion gradients (Na$^+$: $10\,\text{mM}$ inside,
$145\,\text{mM}$ outside; K$^+$: $140\,\text{mM}$ inside,
$4\,\text{mM}$ outside). The pump is the primary active transporter
that funds the cell's resting potential, the action potential's
restoration phase, and a long list of secondary active transporters
that use the inward sodium gradient as their energy source.}
\label{fig:vol06:ch01:na-k-pump}
\end{figure}


\engterm{Worked example 3: the plant cell}. A plant cell is bounded
by a plasma membrane on the inside and a rigid cell wall on the
outside. The wall is built from cellulose microfibrils embedded in a
matrix of hemicellulose and pectin; the microfibrils are oriented in
parallel within a layer and at different orientations between
layers, producing a fibre-reinforced composite that is one of the
strongest natural materials for its density. The wall's tensile
strength is approximately $100\,\text{MPa}$ and its Young's modulus
$10\,\text{GPa}$, comparable to engineering polymers.

The plant cell's central vacuole, occupying $80\,\%$ to $90\,\%$ of
the cell volume in mature plant cells, maintains an osmotic pressure
(turgor) of $0.5$ to $1.0\,\text{MPa}$ against the cell wall. The
turgor pressure is what gives non-woody plants their structural
rigidity: a soft herbaceous stem stands upright because each cell is
pressurised against its wall, summing to a macroscopic load-bearing
capacity. A wilted plant has lost turgor and cannot bear the load;
the cell wall alone is too compliant in shear. The plant cell's
engineering is a pressure-vessel design with the wall as the shell
and the vacuole as the pressuriser.

The pressure-vessel reading allows a direct calculation. A spherical
cell of radius $r$ pressurised to $P$ has a wall tension
$T = Pr/2$ for a thin-shell idealisation. With $r = 20\,\si{\micro\meter}$
and $P = 0.5\,\text{MPa}$, the tension is $T = 5\,\text{N}/\text{m}$,
a value the cellulose wall carries comfortably given its tensile
strength of $100\,\text{MPa}$ over a wall thickness of
$\sim 0.5\,\si{\micro\meter}$. The same calculation applied to a
$1\,\text{cm}$ ripe tomato cell with a thinner wall demonstrates the
edge of the design envelope: split-skin failure under sudden water
uptake is the engineering analogue of dielectric breakdown in a
capacitor.

Chloroplasts are present at $10$ to $100$ per cell in photosynthetic
tissue. Each contains a stacked-membrane system (the thylakoids) that
houses the photosynthetic light reactions, and a soluble compartment
(the stroma) that houses the Calvin cycle. The energy throughput of a
chloroplast and the efficiency of plant photosynthesis are subjects
of Chapter~\ref{vol06:ch02}.

\engterm{Worked example 4: the bacterial cell}. \emph{Escherichia
coli} is the engineer's favourite bacterium because its rapid growth,
small genome, and amenability to genetic manipulation have made it
the workhorse of molecular biology and biotechnology since the 1960s.
A typical \emph{E.~coli} cell is a rod $2\,\si{\micro\meter}$ long
by $0.5\,\si{\micro\meter}$ in diameter, with a doubling time of
$30\,\text{minutes}$ on rich medium at $37\,\degC$. The cell has a
single circular chromosome of $4.6 \times 10^6$ base pairs encoding
$4288$ proteins, no nucleus, no membrane-bounded organelles, and a
cell wall containing peptidoglycan that resists the cell's internal
osmotic pressure of $0.3$ to $0.5\,\text{MPa}$.

The cell's surface is decorated with several types of appendage. The
\engterm{flagella} are helical filaments $20\,\text{nm}$ in diameter
and up to $20\,\si{\micro\meter}$ long, each powered by a rotary
motor embedded in the cell envelope. The motor uses the proton
gradient across the cell membrane as its energy source and rotates at
$100$ to $300\,\text{Hz}$, propelling the cell at $20$ to
$50\,\si{\micro\meter}/\text{s}$. The motor's chemiosmotic coupling
to the proton gradient is the same energy-coupling motif that
Mitchell identified for ATP synthesis in the 1961 chemiosmotic
hypothesis \cite{paper:v6c01-mitchell-1961}; the proton gradient is
the cell's master energy currency, with ATP and rotary motion as two
parallel outputs. The motor is one of the canonical examples of a
rotary molecular machine (Chapter~\ref{vol06:ch04}). Other surface
appendages include \engterm{pili} (fine adhesive filaments) and the
\engterm{sex pilus} that mediates plasmid transfer between cells.

The cell's biochemical throughput, at peak growth on glucose minimal
medium, is approximately $10^{-13}\,\text{g}$ of dry biomass per
cell per generation, requiring approximately $10^{10}$ ATP per
generation. The cell is, by mass and by ATP, the most thermodynamically
efficient producer of biomass per unit substrate that the biotechnology
industry has access to. Chapter~\ref{vol06:ch10} will return to the
economics of bacterial biomanufacturing.

\begin{estimation}
\textbf{Question.} Estimate how many ribosomes a fast-growing
\emph{E.~coli} cell contains at any one moment.

\smallskip
\textit{Estimate, before reading on.}

\smallskip
The total protein content of an \emph{E.~coli} cell at division is
roughly $1.5 \times 10^{-13}\,\text{g}$, or about $10^9$ protein
molecules of average mass $4 \times 10^4\,\text{Da}$. The cell
doubles in $30\,\text{minutes} = 1800\,\text{s}$, so the rate of new
protein synthesis at steady state is $10^9 / 1800 \approx 5.5 \times
10^5\,\text{proteins}/\text{s}$, with each protein on average
$350$ amino acids. The peptide-bond formation rate per ribosome is
approximately $20\,\text{amino acids}/\text{s}$.

The number of ribosomes needed is the amino-acid production rate
divided by the per-ribosome rate:
\[
N_{\text{rib}} \approx \frac{5.5 \times 10^5 \times 350}{20}
\approx 10^7.
\]

Direct measurements give roughly $7 \times 10^4$ ribosomes per
\emph{E.~coli} cell at fastest growth \cite[ch.~17]{text:lehninger-biochemistry}.
The estimate is high by a factor of $100$.

\smallskip
\textit{Postmortem.} The estimate overstates the total protein
synthesis rate because the protein content of a daughter cell is
inherited rather than synthesised from scratch every generation: only
the new biomass needs to be made. Halving the synthesis rate brings
the estimate within a factor of $50$ of the measurement. The
remaining factor traces to two corrections: only $20\,\%$ to $30\,\%$
of the cell's mass is protein at the level relevant to ribosome demand
(the rest is RNA, lipid, polysaccharide, and small molecules), and
the average peptide-bond rate falls when amino acids run short. The
canonical reading of the result is that an \emph{E.~coli} cell at
maximum growth devotes roughly half of its dry mass to ribosomes
themselves; the ribosome is the rate-limiting reagent in bacterial
proliferation. The estimate, even after correction, shows why
biotechnological yields scale with growth rate: the cell's ribosome
content is the throughput bottleneck.
\end{estimation}

\begin{estimation}
\textbf{Question.} Estimate the resting membrane potential of a
mammalian cell from first principles, using only the canonical
intracellular and extracellular potassium concentrations.

\smallskip
\textit{Estimate, before reading on.}

\smallskip
The resting membrane is dominated by potassium permeability: the
plasma membrane has roughly $20$ times the permeability to K$^+$
that it has to Na$^+$. A membrane permeable only to one ion sits at
the Nernst potential for that ion: $V = (RT/zF) \ln([K^+]_o /
[K^+]_i)$. With $R = 8.314\,\text{J}/(\text{mol K})$, $T = 310\,\text{K}$,
$z = +1$, $F = 96485\,\text{C}/\text{mol}$, and concentrations
$[K^+]_o = 4\,\text{mM}$, $[K^+]_i = 140\,\text{mM}$,
\[
V_K = \frac{(8.314)(310)}{(1)(96485)} \ln\left(\frac{4}{140}\right)
\approx (0.0267) \cdot (-3.56)
\approx -0.0951\,\text{V} = -95\,\text{mV}.
\]

The directly measured resting potential of a mammalian neuron is
about $-70\,\text{mV}$. The estimate is too negative by $25\,\text{mV}$.

\smallskip
\textit{Postmortem.} The single-ion estimate ignores the contribution
of Na$^+$ and Cl$^-$. The Goldman-Hodgkin-Katz equation accounts for
all three with their relative permeabilities; with $P_K : P_{Na} :
P_{Cl} = 1 : 0.04 : 0.45$, the predicted resting potential is
$-70\,\text{mV}$, matching the measurement. The Nernst estimate is
the K$^+$ equilibrium potential, and it tells us where the resting
potential \emph{would} sit if the membrane were exclusively
K$^+$-permeable; the offset toward zero is the inward leak through
the small but nonzero Na$^+$ permeability. The reading is engineering
in its purest form: the dominant ion sets the leading order, the
minority ions set the residual, and the controller (the Na$^+$/K$^+$
ATPase) closes the loop by restoring whatever the leak takes away.
The Hodgkin-Huxley framework
\cite{paper:v6c01-hodgkin-huxley-1952,text:v6c01-hille-channels-2001}
is the canonical reference for the full quantitative model.
\end{estimation}

\section{Failure: how cells die}\pathtag{core}

Cells fail by named mechanisms. The classification of cell death has
been refined over the past three decades; for engineering purposes the
three principal modes are apoptosis (programmed death), necrosis
(catastrophic death), and senescence (a stable post-mitotic state
that is not death proper but that is the failure of the cell's
proliferative function). Each has a characteristic mechanism, a
characteristic time-course, and a characteristic consequence for the
surrounding tissue. The 2018 Nomenclature Committee on Cell Death
classification \cite{paper:v6c01-galluzzi-2018} is the current
canonical taxonomy and adds several less common modes (necroptosis,
pyroptosis, ferroptosis, parthanatos, NETosis, lysosome-dependent
death) that we name but do not develop in detail.

\engterm{Apoptosis}. Apoptosis was identified and named as a distinct
mode of cell death by Kerr, Wyllie, and Currie in 1972
\cite{paper:v6c01-kerr-apoptosis-1972,paper:v6c01-kerr-1972-secondary},
who recognised that the morphology (cell shrinkage, chromatin
condensation, membrane blebbing into apoptotic bodies) was a
controlled disassembly rather than a passive collapse. Apoptosis is
the cell's controlled self-disassembly, executed by a conserved
family of cysteine proteases called caspases. The cell shrinks, its
chromatin condenses, its DNA is cleaved into internucleosomal
fragments (the $\sim 180\,\text{bp}$ ladder that is the classical
biochemical signature), and its plasma membrane vesiculates into
membrane-bounded apoptotic bodies that are phagocytosed by
neighbouring cells without releasing pro-inflammatory contents into
the surrounding tissue. The whole process takes a few hours.

The caspase cascade has two principal initiating arms
(\Cref{fig:vol06:ch01:apoptosis-cascade}). The \engterm{intrinsic
(mitochondrial) path} begins with intracellular stress signals (DNA
damage detected by p53, growth-factor withdrawal, hypoxia), shifts
the balance of the Bcl-2 family of regulatory proteins, opens the
mitochondrial outer membrane to release cytochrome $c$ into the
cytoplasm, and activates initiator caspase-9
\cite{paper:v6c01-vaux-korsmeyer-1999}. The \engterm{extrinsic
(death-receptor) path} begins with the binding of external ligands
(Fas ligand, TNF-$\alpha$, TRAIL) to their receptors on the cell
surface, assembles a death-inducing signalling complex (DISC) at the
cytoplasmic face, and activates initiator caspase-8. Both initiator
caspases cleave and activate the executioner caspases (caspase-3 and
caspase-7), which then cleave a few hundred substrate proteins,
including the inhibitor of caspase-activated DNase (releasing the
DNase that cleaves chromatin), structural proteins of the cytoskeleton
and nuclear envelope, and the cell's own repair machinery. The Bcl-2
family controls the threshold: anti-apoptotic Bcl-2 and Bcl-xL keep
the mitochondrial outer membrane intact; pro-apoptotic Bax and Bak
form oligomeric pores that release cytochrome $c$. The cell's
apoptotic commitment is the tipping point at which the pro/anti ratio
crosses a threshold, and the architecture is bistable in the
engineering sense developed in section~1.4.

% Figure: the apoptotic caspase cascade. Two trigger paths (intrinsic /
% mitochondrial; extrinsic / death receptor), converging on the
% executioner caspases and their substrates.
\begin{figure}[htbp]
\centering
\begin{tikzpicture}[
  trig/.style={draw=engred!80!black, fill=engred!15, very thick,
    rectangle, minimum width=2.4cm, minimum height=0.9cm,
    align=center, font=\small, rounded corners=3pt},
  signal/.style={draw=engamber!80!black, fill=engamber!15, very thick,
    rectangle, minimum width=2.4cm, minimum height=0.9cm,
    align=center, font=\small, rounded corners=3pt},
  exec/.style={draw=engblue!80!black, fill=engblue!18, very thick,
    rectangle, minimum width=2.4cm, minimum height=0.9cm,
    align=center, font=\small, rounded corners=3pt},
  outc/.style={draw=engblack, fill=engblue!8, very thick,
    rectangle, minimum width=2.6cm, minimum height=0.9cm,
    align=center, font=\small, rounded corners=3pt},
  arr/.style={-{Latex[length=2.2mm]}, thick, draw=engblack},
  inh/.style={-{Bar[width=3mm]}, thick, draw=engred},
]

% Intrinsic (left) path
\node[trig] (dna) at (0, 4.2) {DNA damage \\ \scriptsize (p53)};
\node[signal] (bcl) at (0, 2.8) {Bcl-2 / Bax \\ balance};
\node[signal] (cyt) at (0, 1.4) {cytochrome $c$ \\ release};
\node[signal] (c9) at (0, 0.0) {caspase-9 \\ (initiator)};

% Extrinsic (right) path
\node[trig] (fas) at (5.5, 4.2) {Fas / TNF-$\alpha$ \\ \scriptsize ligation};
\node[signal] (disc) at (5.5, 2.8) {DISC \\ assembly};
\node[signal] (c8) at (5.5, 0.0) {caspase-8 \\ (initiator)};

% Convergence (centre / bottom)
\node[exec] (c3) at (2.75, -1.6) {caspase-3 / -7 \\ (executioner)};
\node[outc] (mem) at (-0.5, -3.4) {membrane \\ blebbing};
\node[outc] (dna2) at (2.75, -3.4) {DNA \\ fragmentation};
\node[outc] (eat) at (6.0, -3.4) {phagocytosis \\ (no inflammation)};

% Forward arrows: intrinsic
\draw[arr] (dna) -- (bcl);
\draw[arr] (bcl) -- (cyt);
\draw[arr] (cyt) -- (c9);
\draw[arr] (c9) -- (c3);

% Forward arrows: extrinsic
\draw[arr] (fas) -- (disc);
\draw[arr] (disc) -- (c8);
\draw[arr] (c8) -- (c3);

% Executioner -> outcomes
\draw[arr] (c3) -- (mem);
\draw[arr] (c3) -- (dna2);
\draw[arr] (c3) -- (eat);

% Inhibition: Bcl-2 inhibits cytochrome c release (negative regulator
% drawn as bar)
\node[font=\scriptsize, align=center, draw=engred!80!black,
  fill=engred!8, rounded corners=2pt] (bcl2)
  at (-2.8, 2.1) {Bcl-2 \\ (survival)};
\draw[inh] (bcl2) -- (cyt);

% Cross-talk: caspase-8 can also cleave Bid -> mitochondrial path
\draw[arr, densely dashed] (c8.west) to[bend right=20]
  node[midway, above, font=\scriptsize] {Bid} (cyt.east);

\end{tikzpicture}
\caption[The apoptotic caspase cascade]{The two principal triggers
of apoptosis and their convergence on the executioner caspases. The
\textbf{intrinsic} (mitochondrial) path begins with DNA damage
detected by p53, shifts the Bcl-2 / Bax balance, releases cytochrome
$c$ from the mitochondrial inner-membrane space, and activates
caspase-9. The \textbf{extrinsic} (death-receptor) path begins with
Fas ligand or TNF-$\alpha$ binding, assembles the death-inducing
signalling complex (DISC), and activates caspase-8. Both initiators
cleave and activate caspase-3 and caspase-7, the executioner caspases,
which then cleave a few hundred substrate proteins, fragment DNA, and
drive membrane blebbing without releasing pro-inflammatory contents
into the surrounding tissue. Bcl-2 inhibits cytochrome $c$ release
(red bar); caspase-8 also cleaves Bid, linking the extrinsic path to
the mitochondrial path (dashed) \cite{paper:v6c01-vaux-korsmeyer-1999,
paper:v6c01-galluzzi-2018}.}
\label{fig:vol06:ch01:apoptosis-cascade}
\end{figure}


Apoptosis is the body's mechanism for removing damaged, infected, or
supernumerary cells without inflammation. It is essential to
development (the webbed-foot stage in vertebrate limb development is
resolved by apoptosis of the interdigital cells; mammalian breast
involution after lactation is resolved by apoptosis of the milk-
producing epithelium) and to homeostasis (the gut epithelium turns
over every $4$ to $7$ days, balancing proliferation at the crypt with
apoptosis at the villus tip; the immune system maintains its
repertoire by apoptotic deletion of self-reactive lymphocytes). A
human adult turns over roughly $10^{11}$ cells per day by apoptosis,
the same order as new cells generated by mitosis. The failure of
apoptosis is one of the hallmarks of cancer
(Chapter~\ref{vol06:ch03}): tumour cells that have accumulated
oncogenic mutations should have been deleted by apoptotic
surveillance but escape it, typically by upregulating Bcl-2 family
members or by inactivating p53. Cancer therapy that pushes tumour
cells back across the apoptotic threshold is the engineering goal
of much of modern oncology.

\engterm{Necrosis}. Necrosis is the uncontrolled, traumatic death of a
cell. The cell swells, its plasma membrane ruptures, and its
intracellular contents spill into the surrounding tissue, where they
trigger an inflammatory response. The mechanism is the breakdown of
the cell's energy supply: when ATP falls below the level required to
maintain the ion pumps that keep sodium out and potassium in, water
follows the inward sodium gradient and the cell swells to bursting.
The trigger is anything that severely compromises the cell's energy
supply or membrane integrity: ischaemia (loss of blood supply, hence
of oxygen), toxic chemicals that disrupt the mitochondrial respiratory
chain, mechanical damage, heat denaturation. The released contents
include digestive enzymes from lysosomes (which then degrade the
surrounding tissue), DNA and ATP (which act as damage-associated
molecular patterns recognised by the innate immune system,
Chapter~\ref{vol06:ch07}), and cellular debris that the immune system
must clear. Necrosis is the engineering failure mode of the cell
under load: it is what happens when the controlled disequilibrium
fails. The clinical category of infarction (myocardial infarction,
cerebral infarction) is a region of tissue that has undergone necrosis
from ischaemia.

The cascade from ischaemia to necrosis is a textbook chain of
failures. Within seconds of oxygen deprivation, mitochondrial ATP
production falls and the Na$^+$/K$^+$ ATPase rate drops; sodium
accumulates intracellularly and water follows; the cell swells. Within
minutes, lysosomal membranes destabilise as intracellular pH falls,
releasing cathepsins and other proteases; the cytoskeleton begins to
disassemble; mitochondria themselves swell and rupture. By $20$ to
$60\,\text{minutes}$ the plasma membrane fails. Each step has a
characteristic timescale that clinical decision-making (thrombolysis
window, surgical revascularisation) is built around. The 2018 Galluzzi
classification also names a regulated variant, \engterm{necroptosis},
mediated by the RIPK1/RIPK3 kinases that assemble a complex called
the necrosome and recruit MLKL to form pore-like structures in the
plasma membrane \cite{paper:v6c01-galluzzi-2018}; the morphology is
necrosis-like but the trigger is programmed.

\engterm{Senescence}. Senescence is a stable, non-replicative state
that a cell enters in response to certain types of stress. The
senescent cell remains metabolically active, secretes a characteristic
set of cytokines and growth factors (the senescence-associated
secretory phenotype, SASP \cite{paper:v6c01-campisi-2013}), and
resists apoptosis, but it does not divide. Senescence is triggered by
telomere shortening (the \engterm{Hayflick limit}, the empirical
observation that primary cells in culture divide a finite number of
times before stopping; first published by Leonard Hayflick in 1965,
who observed that diploid human fibroblasts divide approximately $50$
times before stopping
\cite{paper:v6c01-hayflick-1965,paper:v6c01-blasco-telomeres-2005}),
by oncogene activation (an unsuccessful attempt at unrestrained
proliferation that triggers protective shutdown), and by persistent
DNA damage. The molecular markers of senescence include induction of
the cyclin-dependent-kinase inhibitor p16INK4a (which blocks G1/S
progression), persistent DNA-damage signalling at telomeres, and
expression of senescence-associated $\beta$-galactosidase
(SA-$\beta$-gal) detectable by histochemical stain at pH $6.0$.
Senescence is engineering's failure-by-saturation mode: the cell's
proliferative capacity has been used up. The accumulation of
senescent cells with age is implicated in many of the chronic diseases
of ageing; the targeted removal of senescent cells (senolytics) is an
active translational area (current as of 2026).

\engterm{Other modes}. Cell-death taxonomy has expanded over the past
twenty years to include several other named modes: \engterm{autophagy}
(a self-eating process that is usually a survival mechanism but can be
lethal at extremes), \engterm{pyroptosis} (an inflammatory cell death
mediated by gasdermin pore formation, used by the innate immune system
to expose infected cells), and \engterm{ferroptosis} (iron-dependent
lipid peroxidation, important in some neurodegenerative contexts). The
proliferation of named modes reflects the biological diversity of
failure rather than a single underlying mechanism. The engineering
reading is that cells die by several distinct mechanisms, each with
its own trigger, time-course, and tissue-level consequence; the
engineer who treats cell death as a single category misses the design
questions that the taxonomy is trying to answer.

\begin{table}[htbp]
\centering
\caption[Three principal cell-death modes compared]{Three principal
cell-death modes compared as engineering failure modes. The full data
file is at
\repopath{volumes/06-life/ch01-cells-structure-function/data/apoptosis\_vs\_necrosis.csv}.}
\label{tab:vol06:ch01:death-modes}
\begin{tabular}{p{3.0cm} p{3.7cm} p{3.7cm} p{3.7cm}}
\hline
\textbf{Feature} & \textbf{Apoptosis} & \textbf{Necrosis} &
  \textbf{Senescence} \\
\hline
Trigger & DNA damage; death-receptor ligation &
  Energy collapse; membrane disruption &
  Telomere shortening; persistent stress \\
ATP requirement & Required throughout & Failed precondition &
  Maintained \\
Timescale & Hours & Minutes to hours & Days to lifetime \\
Cell volume & Shrinks & Swells & Unchanged \\
Membrane integrity & Maintained until late &
  Early rupture & Maintained \\
DNA & Internucleosomal ladder &
  Random fragmentation & Unchanged or stable damage \\
Caspase activity & High (caspase 3/7) & Variable & Absent \\
Inflammation & Absent & Strong (DAMPs) & Mild and chronic (SASP) \\
Clearance & Phagocytosed by neighbours &
  Phagocytosed after immune recruitment &
  Not cleared; accumulates \\
Diagnostic stain & TUNEL; annexin V; active caspase-3 &
  LDH release; propidium iodide & SA-$\beta$-gal \\
First described & 1972 (Kerr et al.) & Pre-modern &
  1961--1965 (Hayflick) \\
\hline
\end{tabular}
\end{table}

\begin{principle}[Cell death is mechanism-specific]
A cell's mode of death determines the consequence for the surrounding
tissue. Apoptosis leaves no inflammatory trace and is the body's
clean-removal mechanism. Necrosis spills cellular contents and
provokes inflammation; the inflammation is the body's wound-response
machinery. Senescence stops the cell from dividing but keeps it alive
and secretory, and the accumulating senescent population is a slow
chronic insult. The engineer designing a therapy, a tissue, or a drug
chooses among the modes: cancer therapy aims to push tumour cells
into apoptosis or pyroptosis; tissue-engineering scaffolds aim to
recruit host cells without inducing necrosis; anti-ageing
interventions aim to reduce senescent burden. Naming the mode is the
first move.
\end{principle}

\section{The cell membrane as engineered interface}\pathtag{standard}

The plasma membrane is, in the cross-domain reading that opens
Volume VI, an instance of the engineered-interface archetype: a thin
boundary that separates two regions with different chemical
composition, supports selective transport between them, maintains the
disequilibrium against passive relaxation, and reports the state of
both sides to internal controllers. Membranes appear in every chapter
of Volume VI in different forms (the mitochondrial inner membrane in
Chapter~\ref{vol06:ch02}, the nuclear envelope in Chapter~\ref{vol06:ch03},
the synaptic membrane in Chapter~\ref{vol06:ch07}), and they
appear in non-biological engineering as well (the membrane in a
reverse-osmosis water purifier, the dielectric in a capacitor, the
ion-exchange membrane in a fuel cell or a chlor-alkali electrolysis
cell). The cross-domain reading is that the same engineering
problem (controlled selectivity across a thin barrier) is being
solved repeatedly by different mechanisms.

\engterm{Selectivity mechanisms}. The plasma membrane is selective by
three distinct mechanisms. \engterm{Lipid solubility}: small,
hydrophobic, uncharged molecules (oxygen, carbon dioxide, ethanol,
steroid hormones) pass directly through the bilayer by diffusion.
\engterm{Channel-mediated passage}: ions and water cross through
protein channels that gate by voltage, ligand binding, mechanical
stress, or temperature. \engterm{Carrier-mediated transport}: larger
or polar molecules (glucose, amino acids) are bound by carrier
proteins that undergo conformational change to translocate the
substrate. Each mechanism is selective in a different way: lipid
solubility is selective by chemistry, channels by geometry and charge,
carriers by molecular recognition. Hille's \emph{Ion Channels of
Excitable Membranes} \cite{text:v6c01-hille-channels-2001} is the
canonical reference for channel biophysics.

\engterm{Active transport}. Many of the membrane's transport tasks are
against the prevailing concentration gradient (sodium out against its
gradient, calcium out against its gradient, hydrogen ion out against
its gradient in many cells) and so require an energy input. The
\engterm{primary active transporters} use ATP directly: the
sodium-potassium ATPase exchanges three sodium ions for two potassium
ions per ATP hydrolysed \cite{paper:v6c01-skou-1957}; the calcium
ATPase pumps two calcium ions per ATP. The \engterm{secondary active
transporters} use an electrochemical gradient established by a
primary transporter as their energy source: the sodium-glucose
cotransporter in the intestinal brush border uses the inward sodium
gradient (which the sodium-potassium ATPase maintains) to pull
glucose against its gradient into the cell. The hierarchy is one of
the volume's clearest illustrations of free-energy coupling: ATP
funds the primary transport, the primary transport funds the
secondary transport, and the secondary transport delivers the
metabolite that, by being catabolised, regenerates the ATP. The
chemiosmotic hypothesis of Mitchell \cite{paper:v6c01-mitchell-1961}
established the same coupling at the mitochondrial inner membrane:
the proton gradient funds ATP synthesis and the bacterial flagellar
motor in parallel.

\engterm{The membrane as a reporting layer}. The membrane carries
signalling receptors that report the extracellular environment to
internal controllers. The receptors are integral membrane proteins
whose extracellular domain binds a specific ligand (a hormone, growth
factor, neurotransmitter, or paracrine signal) and whose intracellular
domain activates a downstream signalling cascade upon binding. The
cascades amplify the signal (a single bound receptor can activate
many G-proteins; each G-protein activates many adenyl cyclases; each
cyclase produces many cAMP molecules; cAMP activates many protein
kinases; the kinases phosphorylate many substrates), so that low
concentrations of extracellular ligand produce large intracellular
responses. The amplification factor is set by the lifetime of each
cascade member; it is one of the engineering tuning knobs by which
the cell sets its responsiveness.

\engterm{Cross-domain pattern}. The reverse-osmosis membrane in a
desalination plant solves the same engineering problem as the
biological membrane: selective transport against a chemical gradient,
maintained at steady state by an external energy input. The polymer
membrane (typically a thin polyamide layer on a porous support) is
selective by size and charge exclusion. The biological membrane is
selective by carrier and channel proteins. Both are thin (the
desalination membrane is approximately $200\,\text{nm}$; the
biological bilayer $4\,\text{nm}$), both require an energy input (the
desalination pump; the cellular ATPase), and both have a maximum
throughput limited by the boundary's surface area. The cross-domain
pattern is that selective transport across a thin boundary is a
generic engineering problem, and that nature and industry have arrived
at convergent answers. Other instances of the same archetype the
reader has met or will meet: the gas-separation membranes of an
air-fractionation plant; the proton-exchange membrane of a fuel cell;
the alveolar epithelium of a vertebrate lung; the gill epithelium of
a fish; the blood-brain barrier; the placental barrier. Each is a
selective interface; each obeys some version of the same accounting.

\section{Project}\pathtag{core}

\begin{project}[Three household microorganism samples]
The reader collects three samples of microorganisms from the home,
inspects each under a low-cost optical microscope, and documents the
structures observed against the chapter's vocabulary.

\textbf{Track}. Household microscopy.
\textbf{Hazard class}: low. The samples are non-pathogenic in
ordinary handling. Wash hands after each session.

\textbf{Equipment}. An optical microscope capable of
$400\times$ magnification (a $40\times$ objective and a $10\times$
eyepiece, USD~$200$ to $400$ new or USD~$50$ to $150$ used);
microscope slides and coverslips; a sample of yeast suspension
(commercial dried baker's yeast rehydrated in warm water); a sample
of yoghurt diluted ten-fold in water (containing live
\emph{Lactobacillus} bacteria); a sample of pond water, aquarium
water, or rainwater from a stagnant puddle (containing protists and
small algae).

\textbf{Procedure}.

\begin{enumerate}[leftmargin=*]
\item For each sample, prepare a wet mount: a drop of the sample on a
  slide with a coverslip lowered carefully to avoid air bubbles.
\item Examine each sample at $100\times$ and $400\times$ magnification.
  Sketch what you see, with a scale bar derived from the field-of-view
  at each magnification.
\item For each organism, identify: cell size and shape; cell wall (yes
  or no, and visible thickness if present); apparent motility; visible
  internal structure (nucleus, chloroplasts, granules); whether the
  organism is prokaryotic or eukaryotic based on the structural
  evidence.
\item For each organism, state which of the chapter's worked-example
  archetypes it most closely resembles and which it does not. Justify
  the choice with at least three structural features.
\end{enumerate}

\textbf{Deliverable}. A five- to eight-page report with three labelled
sketches (one per sample) and a summary table of the four observed
properties (size, motility, structural complexity, classification).
Include a paragraph on at least one structural feature that surprised
you and on the limitations of $400\times$ optical microscopy compared
to electron microscopy (which is needed to resolve organelles smaller
than approximately $1\,\si{\micro\meter}$).

\textbf{Time}. Two to four hours of microscopy and observation; two
to four hours of write-up.
\end{project}

\section{Exercises}

\subsection*{Calculation}\pathtag{core}

\begin{exercise}
A spherical bacterial cell has a radius of $1\,\si{\micro\meter}$.
Compute its volume and its surface area. Compute the surface-to-volume
ratio and compare it to the surface-to-volume ratio of a sphere of
radius $10\,\si{\micro\meter}$ (a typical small animal cell).
\end{exercise}

\paragraph{Solution sketch.}
For a sphere, $V = \tfrac{4}{3}\pi r^3$ and $A = 4\pi r^2$, so
$A/V = 3/r$. At $r = 1\,\si{\micro\meter}$,
$V \approx 4.19\,\si{\micro\meter}^3$,
$A \approx 12.57\,\si{\micro\meter}^2$, and
$A/V = 3\,\si{\micro\meter}^{-1}$. At $r = 10\,\si{\micro\meter}$,
$V \approx 4190\,\si{\micro\meter}^3$,
$A \approx 1257\,\si{\micro\meter}^2$, and
$A/V = 0.3\,\si{\micro\meter}^{-1}$. The bacterial cell has $10\times$
the surface-to-volume ratio of the animal cell. The reading: the
bacterial cell can run all of its transport across the plasma
membrane; the animal cell cannot, and must internalise membrane (ER,
Golgi, mitochondrial cristae) to expose more area
\cite[ch.~1]{text:v6c01-alberts-mbc-2022}. The factor of $10$ in
$A/V$ is the engineering reason for compartmentalisation in
eukaryotes.

\begin{exercise}
A red blood cell modelled as a biconcave disc has a surface area of
$140\,\si{\micro\meter}^2$ and a volume of $90\,\si{\micro\meter}^3$.
Compute its surface-to-volume ratio. Compare to a sphere of the same
volume.
\end{exercise}

\paragraph{Solution sketch.}
The RBC's $A/V = 140/90 \approx 1.56\,\si{\micro\meter}^{-1}$. A
sphere of the same volume has radius
$r = (3V/4\pi)^{1/3}
= (3 \cdot 90 / (4\pi))^{1/3} \approx 2.78\,\si{\micro\meter}$,
giving $A_{\text{sphere}} = 4\pi r^2 \approx 97\,\si{\micro\meter}^2$
and $A/V \approx 1.08\,\si{\micro\meter}^{-1}$. The biconcave disc has
$\sim 1.45\times$ the surface-to-volume ratio of the
equivalent-volume sphere, which is the engineering reason for the
shape: the disc speeds diffusive O$_2$ exchange across the plasma
membrane by $\sim 45\%$ relative to a sphere of the same haemoglobin
content. The deformability of the biconcave shape also lets the cell
squeeze through capillaries of $3\,\si{\micro\meter}$ diameter
\cite[ch.~22]{text:v6c01-alberts-mbc-2022}.

\begin{exercise}
The plasma membrane has a specific capacitance of $0.9\,\si{\micro\farad}/
\text{cm}^2$. A spherical cell of radius $10\,\si{\micro\meter}$
holds a resting potential of $-70\,\text{mV}$. Compute the total
charge separated across the membrane and the number of charges
involved.
\end{exercise}

\paragraph{Solution sketch.}
The cell's surface area is
$A = 4\pi r^2 = 4\pi (10 \times 10^{-4}\,\text{cm})^2
\approx 1.26 \times 10^{-5}\,\text{cm}^2$.
Total capacitance
$C = C_m \cdot A = (0.9 \times 10^{-6}\,\text{F}/\text{cm}^2)
\cdot 1.26 \times 10^{-5}\,\text{cm}^2
\approx 1.13 \times 10^{-11}\,\text{F} = 11.3\,\text{pF}$.
Charge separated $Q = C \cdot V
= 1.13 \times 10^{-11}\,\text{F} \cdot 0.070\,\text{V}
\approx 7.9 \times 10^{-13}\,\text{C}$. The number of charges is
$N = Q/e = 7.9 \times 10^{-13} / (1.6 \times 10^{-19})
\approx 4.9 \times 10^6$ elementary charges. Reading: the charge
separated to hold the resting potential is tiny compared to the
total ion content of the cell ($\sim 10^{11}$ K$^+$ ions in a
cell of $4000\,\si{\micro\meter}^3$ at $140\,\text{mM}$), so the bulk
ion concentrations are essentially unaffected by the potential
itself. The membrane potential is a small-charge phenomenon, not a
bulk-concentration phenomenon.

\begin{exercise}
A sodium-potassium ATPase pumps three sodium ions out and two
potassium ions in per ATP hydrolysed. A neuron at rest hydrolyses
ATP at a rate of $10^7$ molecules per second on this pump. Compute
the ion flux per second and compare to the net ionic charge
crossing during an action potential of duration $1\,\text{ms}$ and
peak inward current $10\,\text{pA}/\si{\micro\meter}^2$ over a
$10\,\si{\micro\meter}^2$ patch.
\end{exercise}

\paragraph{Solution sketch.}
\emph{Pump flux.} At $10^7\,\text{ATP/s}$, the pump moves
$3 \times 10^7\,\text{Na$^+$/s}$ out and $2 \times 10^7\,\text{K$^+$/s}$
in, with net outward current of one positive charge per cycle:
$i_{\text{pump}} = 10^7 \cdot e
\approx 1.6 \times 10^{-12}\,\text{A} = 1.6\,\text{pA}$.
\emph{Action-potential charge.} Current
$I = 10\,\text{pA}/\si{\micro\meter}^2 \cdot 10\,\si{\micro\meter}^2
= 100\,\text{pA} = 10^{-10}\,\text{A}$, over $\Delta t = 10^{-3}\,\text{s}$,
gives charge $Q = I \Delta t = 10^{-13}\,\text{C}$, or
$N = 10^{-13}/1.6 \times 10^{-19} \approx 6.3 \times 10^5$ ions.
\emph{Comparison.} The pump's per-second flux is
$3 \times 10^7 \gg 6.3 \times 10^5$, so the resting pump rate of one
ATP per cycle restores the AP-induced sodium load in a fraction of a
second per patch. At spike rates above a few hundred Hz, however, the
pump becomes the rate-limiting reagent, which is one of the
engineering reasons that neuronal firing rates plateau
\cite{paper:v6c01-attwell-laughlin-2001}.

\begin{exercise}
An \emph{E.~coli} cell of mass $10^{-12}\,\text{g}$ contains $10^7$
ATP molecules at any moment and turns over its entire ATP pool every
second. Compute the cell's ATP throughput in moles per second and per
hour, and convert to a metabolic power in watts.
\end{exercise}

\paragraph{Solution sketch.}
Throughput: $10^7$ ATP/s per cell, divided by Avogadro's number,
gives $n_{\text{ATP}} = 10^7 / 6.022 \times 10^{23}
\approx 1.66 \times 10^{-17}\,\text{mol/s}$, or
$\approx 6 \times 10^{-14}\,\text{mol/hour}$. Free-energy of ATP
hydrolysis under cellular conditions is approximately
$\Delta G \approx -50\,\text{kJ/mol}$. Metabolic power:
$P = n \cdot |\Delta G|
= 1.66 \times 10^{-17}\,\text{mol/s} \cdot 5 \times 10^4\,\text{J/mol}
\approx 8.3 \times 10^{-13}\,\text{W}$. Per gram of \emph{E.~coli}
biomass, this is
$8.3 \times 10^{-13}\,\text{W} / 10^{-12}\,\text{g}
\approx 0.83\,\text{W/g}$, exceeding the specific metabolic rate of a
human ($\sim 0.025\,\text{W/g}$) by $\sim 30\times$. Reading: bacteria
run hot per unit mass; their high metabolic rate is the engineering
driver of their short generation time \cite{text:v6c01-bionumbers-2015}.

\subsection*{Derivation}\pathtag{standard}

\begin{exercise}
Derive the Nernst equation $V_{\text{eq}} = (RT/zF) \ln([\text{ion}]_o/
[\text{ion}]_i)$ for the equilibrium potential of an ion across a
membrane permeable only to that ion. State the assumptions. Apply to
$\text{K}^+$ with $[\text{K}^+]_o = 4\,\text{mM}$ and
$[\text{K}^+]_i = 140\,\text{mM}$ at $310\,\si{\kelvin}$.
\end{exercise}

\paragraph{Solution sketch.}
\emph{Assumptions.} The membrane is selectively permeable to a single
ion species. The system is at electrochemical equilibrium: no net
current flows. The temperature is uniform and the ion activities
equal the concentrations (dilute solution).
\emph{Derivation.} The electrochemical potential of an ion of charge
$z$ at concentration $c$ and electrical potential $\phi$ is
$\tilde\mu = \mu^{\circ} + RT \ln c + zF\phi$. At equilibrium
$\tilde\mu_{\text{in}} = \tilde\mu_{\text{out}}$, so
$RT \ln c_i + zF\phi_i = RT \ln c_o + zF\phi_o$. Rearranging,
$V_{\text{eq}} \equiv \phi_i - \phi_o = (RT/zF)\ln(c_o/c_i)$. The
membrane potential is the value at which the chemical-gradient drive
exactly cancels the electrical-gradient drive.
\emph{Apply to K$^+$.} $R = 8.314\,\text{J}/(\text{mol K})$,
$T = 310\,\text{K}$, $z = +1$, $F = 96485\,\text{C}/\text{mol}$,
$c_o = 4\,\text{mM}$, $c_i = 140\,\text{mM}$. Then
$RT/F = 0.0267\,\text{V} = 26.7\,\text{mV}$, and
$V_K = 0.0267 \cdot \ln(4/140) = 0.0267 \cdot (-3.555)
\approx -95\,\text{mV}$. This is the equilibrium potential for K$^+$
at body temperature with mammalian concentrations; the measured
resting potential ($\sim -70\,\text{mV}$) sits between this value and
the Na$^+$ equilibrium potential because the membrane is K$^+$-dominant
but not K$^+$-only \cite{text:v6c01-hille-channels-2001}.

\begin{exercise}
Derive the Goldman-Hodgkin-Katz equation for the resting membrane
potential of a cell permeable to multiple ion species. Show that it
reduces to the Nernst equation when only one ion is permeable.
\end{exercise}

\paragraph{Solution sketch.}
Assume the membrane is permeable to K$^+$, Na$^+$, and Cl$^-$ with
permeabilities $P_K$, $P_{Na}$, $P_{Cl}$. The Goldman constant-field
assumption is that the electric field is uniform across the membrane
of thickness $d$. The Nernst-Planck flux for each ion is
$J_i = -P_i z_i (FV/RT) [(c_{i,o} - c_{i,i} e^{z_i FV/RT})/(1 - e^{z_i FV/RT})]$.
At steady state and with no net current,
$\sum z_i J_i = 0$. Solving for V (after some algebra) gives the GHK
equation
\[
V = \frac{RT}{F} \ln
\frac{P_K [K^+]_o + P_{Na}[Na^+]_o + P_{Cl}[Cl^-]_i}
     {P_K [K^+]_i + P_{Na}[Na^+]_i + P_{Cl}[Cl^-]_o}.
\]
Cl$^-$ enters with reversed concentrations because of its negative
charge. \emph{Reduction.} If $P_{Na} = P_{Cl} = 0$, the numerator
becomes $P_K [K^+]_o$ and the denominator $P_K [K^+]_i$; the $P_K$
cancels, leaving $V = (RT/F)\ln([K^+]_o/[K^+]_i)$, the Nernst equation
for K$^+$. \emph{Numerical check.} With $P_K : P_{Na} : P_{Cl} =
1 : 0.04 : 0.45$ and the canonical mammalian concentrations,
$V_{\text{rest}} \approx -70\,\text{mV}$, matching measured resting
potentials in many cell types \cite{text:v6c01-hille-channels-2001}.
The code \texttt{membrane\_voltage\_nernst.py} verifies both the
GHK calculation and the Nernst reduction numerically.

\begin{exercise}
A vesicle of radius $r$ has bending energy
$E = 8\pi \kappa$ (independent of $r$ for a closed sphere), where
$\kappa$ is the membrane bending modulus. Use the Helfrich form
$E = \frac{1}{2}\kappa \int (2H - H_0)^2 \,dA$ with $H$ the mean
curvature, $H_0$ the spontaneous curvature, and integration over the
membrane area. Identify the conditions under which a tube and a
sphere have equal bending energy.
\end{exercise}

\paragraph{Solution sketch.}
\emph{Sphere.} For a sphere of radius $r$, $H = 1/r$ everywhere, and
the area is $A = 4\pi r^2$. Setting $H_0 = 0$ (symmetric bilayer) and
using $(2H)^2 = 4/r^2$:
$E_{\text{sph}} = \tfrac{1}{2}\kappa \cdot (4/r^2) \cdot 4\pi r^2
= 8\pi\kappa$, independent of $r$. This is the textbook Helfrich
result \cite{paper:v6c01-helfrich-1973}.
\emph{Tube.} A cylindrical tube of radius $R$ and length $L$ has
$H = 1/(2R)$ (one principal curvature $1/R$, the other zero), so
$(2H)^2 = 1/R^2$ and $A = 2\pi R L$:
$E_{\text{tube}} = \tfrac{1}{2}\kappa \cdot (1/R^2) \cdot 2\pi R L
= \pi\kappa L / R$.
\emph{Equal-energy condition.} $E_{\text{sph}} = E_{\text{tube}}$
gives $8\pi\kappa = \pi\kappa L/R$, or $L = 8R$. A tube of length
$8R$ has the same bending energy as a sphere of any radius. With
$\kappa \approx 10\,k_B T$, both energies are $\sim 80\,k_B T$ for a
small vesicle; this sets the thermal-fluctuation scale for vesicle
shape changes \cite{text:v6c01-phillips-physical-2012}.

\subsection*{Estimation}\pathtag{core}

\begin{exercise}
Estimate the number of mitochondria in a single cardiomyocyte of
volume $20\,000\,\si{\micro\meter}^3$, given that mitochondria
occupy $35\,\%$ of the cell volume and each mitochondrion is
approximately $1\,\si{\micro\meter}^3$.
\end{exercise}

\paragraph{Solution sketch.}
Mitochondrial volume per cell:
$V_{\text{mito}} = 0.35 \cdot 20\,000\,\si{\micro\meter}^3
= 7000\,\si{\micro\meter}^3$. With each mitochondrion occupying
$\sim 1\,\si{\micro\meter}^3$, the count is $N \approx 7000$
mitochondria per cardiomyocyte. The canonical literature value is
approximately $5000$ per cardiomyocyte \cite{text:v6c01-bionumbers-2015};
the factor-of-$1.5$ difference traces to the mitochondrial volume
being $\sim 1.4\,\si{\micro\meter}^3$ on average in adult human
cardiomyocytes rather than $1.0\,\si{\micro\meter}^3$, and to the
$35\%$ volume fraction being an upper bound. The reading: a
cardiomyocyte invests over a third of its volume in ATP-producing
machinery, the highest among any human cell type, which is the
engineering reason that cardiac tissue is the body's most ischaemia-
vulnerable organ.

\begin{exercise}
Estimate the number of cells in an adult human body, given that the
typical cell volume is $1000$ to $10\,000\,\si{\micro\meter}^3$ and
the body's total volume is approximately $70\,\text{L}$. Account for
the fact that erythrocytes are smaller than typical and that
adipocytes are larger.
\end{exercise}

\paragraph{Solution sketch.}
\emph{First-cut estimate.} $70\,\text{L} = 7 \times 10^{13}\,\si{\micro\meter}^3$,
divided by a mean cell volume of $\sim 5000\,\si{\micro\meter}^3$,
gives $\sim 1.4 \times 10^{10}$ cells. But this is wrong by orders of
magnitude because erythrocytes are tiny ($90\,\si{\micro\meter}^3$)
and dominate the count.
\emph{Better accounting.} Bianconi et al.\ 2013 \cite{paper:v6c01-bianconi-2013}
report a careful organ-by-organ accounting that totals
$3.7 \times 10^{13}$ cells, of which $\sim 2.5 \times 10^{13}$ are
erythrocytes (the rest split among muscle, neurons, adipocytes,
epithelium, and endothelium). The number $3.7 \times 10^{13}$
supersedes the older $10^{13}$ to $10^{14}$ range. Bacterial cells
associated with the host are now estimated at $\sim 3.8 \times 10^{13}$
\cite{paper:v6c01-sender-2016}, putting the human-to-microbiota ratio
near $1:1$ rather than the older 1:10 figure. Reading: estimation by
volume division is fast but blind to compositional heterogeneity;
the careful enumeration is order-of-magnitude consistent with the
volume estimate but moves the answer by a factor of $3$.

\begin{exercise}
Estimate the rate of glucose consumption by the human brain in
$\text{g}/\text{day}$, given that the brain consumes about
$20\,\%$ of the body's $2000\,\text{kcal}/\text{day}$ at rest and
that glucose has an energy content of $4\,\text{kcal}/\text{g}$ when
fully oxidised.
\end{exercise}

\paragraph{Solution sketch.}
Brain energy budget: $0.20 \cdot 2000\,\text{kcal/d}
= 400\,\text{kcal/d}$. Glucose required:
$400\,\text{kcal/d} / 4\,\text{kcal/g} = 100\,\text{g/d}$ of glucose.
This matches measured cerebral glucose uptake (approximately
$120\,\text{g/d}$ in adults, slightly higher because the brain also
uses lactate and ketone bodies as alternative fuels)
\cite{paper:v6c01-erecinska-1989}. With glucose carrying
$\sim 30$ ATP per molecule oxidised, the brain's ATP turnover is
$\sim 6 \times 10^{21}$ ATP/s, dominated by the Na$^+$/K$^+$ ATPase
restoring ion gradients after each action potential
\cite{paper:v6c01-attwell-laughlin-2001}. Reading: a fasting adult
delivering glucose at $\sim 100\,\text{g/d}$ to the brain is
running at the lower bound of liver gluconeogenesis; sustained
hypoglycaemia below this threshold produces the classic neurological
symptoms of brain energy starvation.

\begin{exercise}
Estimate the diffusion time for a small molecule (diffusion
coefficient $D \sim 10^{-9}\,\text{m}^2/\text{s}$) to traverse a
spherical bacterial cell of radius $1\,\si{\micro\meter}$ and a
spherical animal cell of radius $10\,\si{\micro\meter}$.
\end{exercise}

\paragraph{Solution sketch.}
The mean-square displacement for 3D Fickian diffusion is
$\langle r^2 \rangle = 6Dt$, so the characteristic time to cross
distance $L$ is $\tau = L^2/(6D)$.
\emph{Bacterium.} $L = 10^{-6}\,\text{m}$,
$\tau = 10^{-12}/(6 \times 10^{-9}) \approx 1.7 \times 10^{-4}\,\text{s}
\approx 0.17\,\text{ms}$.
\emph{Animal cell.} $L = 10^{-5}\,\text{m}$,
$\tau = 10^{-10}/(6 \times 10^{-9}) \approx 17\,\text{ms}$.
Crowding-corrected (Luby-Phelps factor of $\sim 3.5$
\cite{paper:v6c01-luby-phelps-2000}), both times stretch by a factor
of $\sim 3.5$ in cytoplasm: $\sim 0.6\,\text{ms}$ and $\sim 60\,\text{ms}$.
Reading: diffusion is fast at the bacterial scale and adequate but
not fast at the animal-cell scale; large cells (axons, oocytes) need
motor-driven transport on the cytoskeleton because diffusion is too
slow on their lengthscale. The code
\texttt{diffusion\_timescale.py} tabulates the result for a range
of $L$ and $D$.

\subsection*{Simulation}\pathtag{mastery}

\begin{exercise}
Implement a simple model of the bacterial chemotaxis adaptation
circuit as a system of two ODEs with one fast variable (receptor
activity) and one slow variable (receptor methylation). Show
numerically that the steady-state receptor activity returns to a
fixed baseline regardless of step-input magnitude.
\end{exercise}

\paragraph{Solution sketch.}
A minimal adaptation model has fast receptor activity $a$ and slow
methylation $m$:
$\dot a = (1/\tau_a)(F(L, m) - a)$,
$\dot m = k_R (1 - a) - k_B a$, where $L$ is the attractant
concentration, $F$ a sigmoidal activation function, and $\tau_a \ll
\tau_m$. At steady state $\dot m = 0 \Rightarrow a_{\text{ss}} =
k_R/(k_R + k_B)$, independent of $L$: integral feedback delivers
exact adaptation \cite{paper:v6c01-yi-2000}. The reader implements
the system in Python (Euler or RK4 with $\Delta t = 0.01$) starting at
$L = 1$, stepping $L$ to $10$ at $t = 50$, and observing $a$ spike
upward then return to the baseline $a_{\text{ss}}$ within $\sim
5\tau_m$. Repeat with $L$ stepping to $100$; the baseline-return is
the same. A scatter plot of (step magnitude, post-step steady state)
shows a flat line at $a_{\text{ss}}$, the signature of exact
adaptation. The mechanism is the algebraic structure: $\dot m$
proportional to $(1-a)$ minus a constant fraction of $a$ makes $m$
integrate the error $(a - a_{\text{ss}})$ over time, which is the
defining feature of an integral controller.

\begin{exercise}
Simulate the coherent type-1 feed-forward loop as a system of three
coupled ODEs with Hill-function regulation. Plot the response of $Z$
to a step-up of $X$ and to a step-down. Quantify the delay and the
filter cutoff.
\end{exercise}

\paragraph{Solution sketch.}
Equations: $X = X_{\text{input}}(t)$ (the driving signal),
$\dot Y = \beta_Y H(X) - \alpha_Y Y$,
$\dot Z = \beta_Z H(X) H(Y) - \alpha_Z Z$, where the AND-gate at $Z$
uses a product of Hill functions $H(u) = u^n/(K^n + u^n)$
\cite[ch.~4]{text:alon-systems-biology,paper:v6c01-shoval-alon-2010}.
Choose $\beta = \alpha = 1$, $K = 0.5$, $n = 4$, $Y_0 = Z_0 = 0$,
and run with $X$ stepping from $0$ to $1$ at $t = 10$. $Y$ rises with
time constant $1/\alpha_Y$; $Z$ stays near zero until $Y$ exceeds
$K$, then rises sharply: this is the delayed step-up response, with
delay $\sim \ln 2/\alpha_Y \approx 0.7$ time units. Now step $X$ back
to $0$ at $t = 60$; $Z$ drops promptly because $H(X) = 0$ makes the
production term vanish immediately, regardless of $Y$: this is the
prompt step-down response. The sign-sensitive delay filters brief
pulses in $X$: a pulse shorter than the $Y$-rise time never crosses
the $K$ threshold at $Y$, so $Z$ stays low. Pulses longer than
$\sim 1/\alpha_Y$ propagate. The cutoff is approximately
$f_c \sim \alpha_Y/(2\pi)$.

\subsection*{Design}\pathtag{standard}

\begin{exercise}
Design a thought-experiment cell that secretes a target molecule at a
controlled rate of $10^4$ molecules per second in response to a
specific extracellular signal. Specify the receptor, the signalling
cascade, the synthesis enzyme, and the feedback mechanism that
maintains the secretion rate against substrate depletion. Identify
the rate-limiting step.
\end{exercise}

\paragraph{Discussion.}
A well-formed answer specifies all five components: \emph{receptor}
(e.g.\ G-protein-coupled receptor binding the signal at the surface,
with a $K_d$ matched to expected ligand range); \emph{signalling
cascade} (G$_s$-protein activates adenylyl cyclase, generating cAMP,
which activates protein kinase A); \emph{synthesis enzyme} (PKA
phosphorylates and activates the rate-limiting enzyme in the
synthesis pathway; identify a candidate substrate flow into the
product); \emph{secretion mechanism} (constitutive or regulated
exocytosis; vesicle fusion at SNARE complexes); \emph{feedback}
(product binds back to a repressor of the synthesis-enzyme gene, an
end-product inhibition loop). The rate-limiting step at the
$10^4\,\text{s}^{-1}$ secretion rate is most likely the vesicle
fusion machinery ($\sim 10\,\text{ms}$ per vesicle, capping single-
site secretion at $\sim 100\,\text{s}^{-1}$ per fusion site, so the
cell needs at least $\sim 100$ active fusion sites operating in
parallel). The rubric values: clear specification of each component;
quantitative match between requested rate and component capacity;
identification of the bottleneck; recognition that substrate
depletion can stall the system without a feedback that adjusts
upstream synthesis.

\begin{exercise}
A vendor proposes a ``synthetic plant cell'' that uses an engineered
plasma membrane to support photosynthesis without a cell wall. Identify
the engineering constraints the design must satisfy (turgor pressure,
membrane stress, osmotic balance) and the design choices it would have
to make to satisfy them.
\end{exercise}

\paragraph{Discussion.}
A well-formed answer identifies the engineering tension: a plant cell
maintains $0.5$ to $1.0\,\text{MPa}$ turgor against its cell wall;
without a wall, the same osmotic pressure would burst the plasma
membrane (whose tensile strength is approximately $5\,\text{mN/m}$
hoop tension before lysis). The vendor must therefore either
(a)~match the membrane's external osmolarity to its internal
osmolarity, eliminating the turgor: but then the cell has no
internal pressure to drive cell expansion, no shape control, and no
mechanical advantage in tissue assembly; (b)~stiffen the membrane
with a synthetic surface coating that bears the pressure: but this
loses the membrane's self-healing and selectivity; (c)~suspend the
cells in a hyperosmotic medium that matches whatever turgor the
chloroplasts require for thylakoid function: but this restricts the
operating environment and complicates harvesting. A high-marking
answer notices that the cell-wall question is connected to
photosynthesis: the thylakoid membrane inside the chloroplast also
maintains a proton gradient and an osmotic pressure, and without
the cell wall to constrain the cell's overall volume, the thylakoid
gradient and the cell-level osmolarity must be balanced independently.
The rubric values: identification of the pressure-vessel role of the
wall; quantitative estimate of the membrane stress without a wall;
discussion of at least two design alternatives with their tradeoffs.

\subsection*{Diagnosis and reverse engineering}\pathtag{core}

\begin{exercise}
A patient's red blood cells appear normal in size and shape under a
microscope but rupture (haemolyse) when placed in distilled water and
shrink when placed in $1\,\text{M}$ saline. Identify the property
under test and the cellular mechanism that produces the observed
behaviour.
\end{exercise}

\paragraph{Solution sketch.}
The property under test is \emph{osmotic fragility}. The patient's
RBCs are behaving as expected for any cell with a semi-permeable
membrane: distilled water (hypotonic, $\sim 0\,\text{mOsm}$) drives
water inward by osmosis, the cell swells, and the plasma membrane
ruptures when the cell's spherical volume is reached at
$\sim 165\,\si{\micro\meter}^3$ (the maximum the membrane can
enclose without lysis); $1\,\text{M}$ saline ($\sim 2000\,\text{mOsm}$,
hypertonic to the cell's $\sim 300\,\text{mOsm}$ cytoplasm) drives
water outward and the cell shrinks (``crenates,'' with a spiky
contour from the cortical cytoskeleton). The cellular mechanism is
the same in both cases: water moves across the bilayer down its
chemical potential gradient through aquaporin channels and through
the lipid bilayer directly. The test is informative because a cell
with abnormally fragile membrane (hereditary spherocytosis, due to
defects in spectrin or band-3 protein) lyses at much higher
osmolarity than a normal RBC; a cell with abnormally resistant
membrane lyses only at lower osmolarity. The reading: even ``normal''
behaviour is diagnostic if the test conditions are calibrated. The
quantitative version of this test is the \emph{osmotic fragility
curve}, plotting fraction haemolysed against medium osmolarity.

\begin{exercise}
A bacterial culture stops growing exponentially and enters stationary
phase after $8\,\text{hours}$. The medium contains an excess of
glucose and oxygen. Identify three possible rate-limiting factors and
propose a measurement that would distinguish among them.
\end{exercise}

\paragraph{Solution sketch.}
Three candidate rate-limiting factors:
(1)~\emph{Nitrogen or phosphorus depletion}: bacteria need both for
protein and nucleic-acid synthesis. Distinguish by measuring residual
ammonium / nitrate / phosphate in the medium; add it back to a
sub-culture and check whether growth resumes.
(2)~\emph{pH excursion}: bacterial fermentation can acidify the
medium below the strain's tolerance ($\sim$ pH 5 for \emph{E.~coli},
$\sim$ pH 4 for \emph{Lactobacillus}). Distinguish by measuring
medium pH at the transition; buffer or adjust to neutral and check
for resumption.
(3)~\emph{Waste product accumulation}: acetate, ethanol, or other
fermentation products inhibit growth above threshold concentrations.
Distinguish by quantifying these waste products (HPLC or specific
enzymatic assay); a dilution into fresh medium that doesn't restore
growth (when nutrients are confirmed) implicates a metabolic state
change rather than a depletion. A fourth possibility, quorum sensing
(autoinducer accumulation triggering programmed stationary-phase
gene expression), can be tested by adding spent supernatant from a
stationary culture to a fresh exponential culture and checking
whether the addition itself induces stationary phase. The rubric
values: identification of distinct mechanism classes; design of a
measurement that distinguishes (not just identifies) candidates.

\begin{exercise}
A skin biopsy reveals a population of cells that are large, flat,
non-replicating, and secreting elevated levels of inflammatory
cytokines. Identify the cellular state and propose a histological
stain that would confirm the diagnosis.
\end{exercise}

\paragraph{Solution sketch.}
The cellular state is \emph{senescence} (likely replicative
senescence with the senescence-associated secretory phenotype, SASP)
\cite{paper:v6c01-campisi-2013}. The morphological hallmarks
(enlarged, flattened cells) and the secretory phenotype (cytokines
including IL-6, IL-8, MMPs) are diagnostic. The standard confirmatory
histological stain is \emph{SA-$\beta$-galactosidase} (senescence-
associated $\beta$-galactosidase), which is detected by X-gal
substrate at pH $6.0$; senescent cells show distinctive blue staining
in lysosomes, distinguishable from non-senescent cells because
non-senescent lysosomal $\beta$-gal activity is too low at pH $6.0$
to give visible staining. Other confirmatory markers include
p16INK4a immunohistochemistry (upregulated in senescent cells),
$\gamma$-H2AX foci (persistent DNA damage), and loss of nuclear
lamin-B1. The reading: the dermal accumulation of senescent
fibroblasts is one of the engineering mechanisms of skin ageing;
removing these cells with senolytics is an active intervention
strategy.

\subsection*{Failure analysis}\pathtag{core}

\begin{exercise}
A region of cardiac tissue shows necrotic histology after a
$30\,\text{minute}$ interruption of its blood supply. Identify the
chain from ischaemia to necrosis at the cellular level (which
membrane fails first, which gradient collapses, which cellular
contents leak, what the inflammatory response will be).
\end{exercise}

\paragraph{Solution sketch.}
The cascade from ischaemia to necrosis:
\emph{Step 1 (seconds to minutes).} Oxygen supply falls, mitochondrial
respiratory chain stops, ATP production collapses, ADP and AMP
accumulate. ATP-dependent processes begin to fail.
\emph{Step 2 (minutes).} The Na$^+$/K$^+$ ATPase rate falls below
the leak rate, intracellular Na$^+$ rises, water follows osmotically,
the cell swells. Intracellular Ca$^{2+}$ rises because the Ca-ATPase
similarly fails and Na$^+$/Ca$^{2+}$ exchanger reverses direction. The
elevated Ca$^{2+}$ activates phospholipases that degrade membrane
lipids and proteases (calpains) that degrade cytoskeleton.
\emph{Step 3 (10--30 minutes).} The lysosomal membrane destabilises
as cytoplasmic pH falls; lysosomal proteases (cathepsins) are
released into the cytoplasm. The mitochondrial outer membrane
permeabilises and releases cytochrome $c$, but the parallel ATP
collapse prevents apoptotic execution; the cell defaults to necrosis
rather than apoptosis. \emph{Step 4 (30--60 minutes).} The plasma
membrane ruptures (the wall of the failed osmotic-pressure vessel),
spilling intracellular contents (LDH, troponin, ATP, DNA, cytochrome
$c$) into the extracellular space. These contents are damage-
associated molecular patterns (DAMPs) recognised by neighbouring
cells' Toll-like receptors and other innate-immune sensors.
\emph{Step 5 (hours to days).} Neutrophils, then macrophages, infiltrate
the necrotic zone and clear debris; cytokine release drives systemic
inflammation; healing proceeds by fibrotic scarring rather than
restoration of contractile tissue (the heart's cardiomyocytes are
post-mitotic). The clinical signature includes troponin elevation
(spilled from the cells; the biomarker of myocardial infarction).

\begin{exercise}
A drug intended to kill cancer cells by triggering apoptosis instead
provokes a strong inflammatory response in the patient. Identify
which step in the apoptotic pathway is likely failing and propose
one diagnostic that would confirm whether the cells are dying by
necrosis or by pyroptosis instead.
\end{exercise}

\paragraph{Solution sketch.}
Three plausible failures:
(1)~\emph{Apoptotic ATP collapse}: the cancer cells may have impaired
mitochondrial function (Warburg metabolism), and the apoptotic
program (which requires ATP) cannot complete. The default fallback
is necrosis.
(2)~\emph{Membrane permeabilisation by the drug}: the drug may be
directly puncturing membranes (or activating gasdermin pore
formation), producing pyroptosis rather than apoptosis. Pyroptosis
releases IL-1$\beta$ and IL-18, both potent inflammatory cytokines.
(3)~\emph{Failure of phagocytic clearance}: even properly executed
apoptosis is inflammatory if apoptotic bodies are not cleared promptly
(secondary necrosis). This is common in solid tumours with poor
vascular access.
\emph{Diagnostic.} The cleanest distinction is measuring released
cytokines: pyroptosis releases IL-1$\beta$ (a hallmark) and active
caspase-1; necrosis releases HMGB1 (a DAMP, but caspase-1 negative).
A second diagnostic is the morphology: pyroptotic cells form
gasdermin-D pores and balloon outward through membrane bulges
distinct from necrotic swelling; necrotic cells swell more uniformly.
The 2018 Galluzzi classification \cite{paper:v6c01-galluzzi-2018}
gives the full battery of distinguishing markers. The therapeutic
implication: a drug that triggers pyroptosis in tumour cells may have
value in cold tumours (the inflammation can recruit immune cells),
but in this patient the response is excessive and the drug dose or
delivery scheme needs revision.

\subsection*{Judgment}\pathtag{standard}

\begin{exercise}
A press release announces a ``self-powering cell'' that draws energy
from a single, undisclosed source and maintains its disequilibrium
indefinitely. Write a one-paragraph reply identifying the
thermodynamic claim being made and the kind of measurement that
would distinguish a genuine breakthrough from a misreading of the
cell's energy supply.
\end{exercise}

\paragraph{Discussion.}
A well-formed reply makes three points. First, the claim of an
``undisclosed'' energy source is the warning sign: living cells obey
the second law of thermodynamics, so maintaining disequilibrium
requires a free-energy input from somewhere. The question is only
where it comes from. Second, the measurement: a calorimetric energy
balance, summing all known inputs (e.g.\ glucose, oxygen, ATP from
medium, photons) and outputs (heat, work, waste), should close to
within experimental error. If it closes, the ``self-powering'' claim
is a misreading of an ordinary cell whose energy input was being
delivered in some unmeasured form (radiation, dissolved nutrient, a
membrane potential maintained by some unnoticed gradient). If the
balance does not close, the cell is either drawing on an undisclosed
input (and the press release is hiding it) or violating the second
law (in which case the experiment is wrong). Third, the relevant
benchmark: known autotrophs (algae, cyanobacteria, lithoautotrophic
bacteria) all close their energy balances at light, CO$_2$, or
chemical reductant. The high-marking answer states that the burden
of proof for the extraordinary claim of new thermodynamics is on the
claimant, and a closing calorimetric balance is the standard test.

\begin{exercise}
A textbook describes the cell as a ``factory''. Identify two ways in
which the factory metaphor is accurate (input/output, regulated
throughput) and two ways in which it is misleading (self-replication
of the factory itself; the absence of an external designer choosing
the throughput).
\end{exercise}

\paragraph{Discussion.}
The factory metaphor is useful and limited. \emph{Accurate.}
(i)~Cells process inputs (glucose, oxygen, amino acids) into outputs
(proteins, secreted molecules, daughter cells) at controlled rates,
matching the factory's input-output schema. (ii)~The throughput is
regulated by feedback loops with well-defined setpoints, mirroring
factory production-line control. \emph{Misleading.}
(iii)~A factory does not build copies of itself; cells do, and the
self-replication is not an add-on but the primary engineering feature.
The closest factory analogue is a self-replicating fabricator, which
is theoretical and rare; cells are the only working examples. (iv)~A
factory's throughput, product specifications, and architecture are
chosen by external designers; a cell's are chosen by evolutionary
selection acting on its replication rate and survival. The cell does
not have a customer; it has a fitness function, which is a different
control objective. A high-marking answer also notices that the
factory metaphor undersells the cell's information-processing role:
cells compute, in the sense of \cite{paper:v6c01-bray-1995}, while
factories simply execute pre-set instructions. The rubric values:
recognition of both lanes; specific examples in each; identification
of self-replication as the structural disanalogy.

\begin{exercise}
A vendor advertises a tissue scaffold ``compatible with all cell
types''. Identify the cellular-level features (membrane composition,
surface chemistry, mechanical compliance, attachment ligands) that
determine compatibility, and write a one-paragraph reply asking the
vendor to specify which of those features they have characterised.
\end{exercise}

\paragraph{Discussion.}
``Compatible with all cell types'' is an unfalsifiable marketing
claim; cells of different types attach, spread, and survive under
quite different conditions. The well-formed reply asks the vendor to
specify:
(1)~\emph{Surface chemistry}: which integrin-binding motifs are
displayed (RGD, fibronectin domain III, laminin LG domains, collagen)
and at what density. Different cell types require different ligands;
a scaffold tuned for endothelial cells will not necessarily attract
osteoblasts.
(2)~\emph{Mechanical compliance}: scaffold elastic modulus matters,
because cells sense substrate stiffness and differentiate accordingly
(soft substrates favour neural differentiation of mesenchymal stem
cells; stiff substrates favour bone).
(3)~\emph{Porosity and topography}: feature size at the
$\sim 10\,\si{\micro\meter}$ scale (cell size) and surface texture
at the nanoscale modulate adhesion and migration.
(4)~\emph{Degradability and degradation products}: the scaffold's
breakdown rate must match the host cells' tissue-formation rate, and
the breakdown products must not be cytotoxic.
(5)~\emph{Sterility and endotoxin content}: any bacterial cell-wall
fragment triggers an inflammatory response that overrides the
intended compatibility. The high-marking reply also asks: which cell
types have been tested on the scaffold, at what cell density, for
how long, and against what controls. The rubric values: identification
of multiple cell-attachment determinants; concrete specification of
what the vendor should report.

\begin{exercise}
A colleague asserts that mitochondria are vestigial organelles whose
remaining function is purely energetic. Identify two non-energetic
roles that mitochondria perform in modern eukaryotic cells (apoptotic
signalling, calcium buffering, biosynthesis of haem and steroid
precursors) and write a one-paragraph reply naming the consequence
that follows for any therapy that disrupts mitochondrial function
without distinguishing among those roles.
\end{exercise}

\paragraph{Discussion.}
The colleague's claim is incorrect on two fronts. \emph{Non-energetic
roles}: mitochondria release cytochrome $c$ to initiate the intrinsic
apoptotic pathway \cite{paper:v6c01-vaux-korsmeyer-1999}; they
sequester intracellular Ca$^{2+}$ during signalling events (the
mitochondrial Ca$^{2+}$ uniporter takes up Ca$^{2+}$ at the
$\sim$ micromolar local concentrations achieved at ER-mitochondria
contact sites); they conduct critical biosynthesis steps (the haem
biosynthesis pathway begins and ends in mitochondria; the synthesis
of steroid hormones starts there; iron-sulfur cluster assembly
occurs there; lipid metabolism uses the matrix as its hub). A drug
that disrupts mitochondrial function without distinguishing among
these roles risks an unintended cascade: mitochondrial dysfunction
that triggers apoptosis (potentially helpful in cancer, dangerous in
heart and brain), that dysregulates Ca$^{2+}$ signalling (cardiac
arrhythmias, neural excitotoxicity), that depletes haem (anaemia),
and that depletes steroid hormones (endocrine effects). The
therapeutic implication is that mitochondrially targeted drugs
require tissue-specific delivery or molecularly specific targeting
of the energetic apparatus (electron-transport-chain complexes, ATP
synthase) without engaging the apoptotic or Ca$^{2+}$ functions; the
endosymbiotic origin \cite{paper:v6c01-margulis-1967} is the
historical reason the organelle has accumulated so many roles, and
the engineering reason for its present complexity. The rubric values:
identification of at least two non-energetic roles; quantitative or
mechanistic linkage between disruption and side effect; recognition
that ``mitochondrial dysfunction'' is too coarse a category to design
against.
